\documentclass[11pt,reqno]{amsart}

\usepackage[T1]{fontenc}
\usepackage{amsmath,amssymb,amsthm}
\usepackage{booktabs}
\usepackage[expansion=false]{microtype}
\usepackage{xcolor}
\usepackage[margin=1.2in]{geometry}
\usepackage[colorlinks=true,linkcolor=blue!50!black,citecolor=blue!50!black,
            urlcolor=blue!50!black]{hyperref}

\theoremstyle{plain}
\newtheorem{theorem}{Theorem}[section]
\newtheorem{proposition}[theorem]{Proposition}
\newtheorem{lemma}[theorem]{Lemma}
\newtheorem{corollary}[theorem]{Corollary}
\theoremstyle{definition}
\newtheorem{remark}[theorem]{Remark}
\newtheorem{hypothesis}{Hypothesis}
\renewcommand{\thehypothesis}{H($\vartheta$)}

\DeclareMathOperator{\Res}{Res}
\newcommand{\R}{\mathbb{R}}
\newcommand{\C}{\mathbb{C}}
\newcommand{\Z}{\mathbb{Z}}
\newcommand{\Rp}{\operatorname{Re}}
\newcommand{\Ip}{\operatorname{Im}}
\newcommand{\lxi}[1]{\left(\tfrac{\xi'}{\xi}\right)\!#1}

\title[The margin in the shifted screw-function criterion]
      {The margin in the shifted screw-function criterion\\ for zero-free half-planes}

% ---------------------------------------------------------------------------
\author{Edward B. Baker III}
\email{edwardbaker86@gmail.com}

\date{\today}

\begin{document}

\begin{abstract}
Suzuki attached to the Riemann zeta-function a continuous real function $\Psi$ --- the screw
function of $\zeta$ --- and proved that the Riemann Hypothesis is equivalent to the pointwise
inequality $\Psi(t)\ge 0$. He also introduced a shifted family $\Psi_\omega$ and proved that
$\zeta(s)\neq 0$ for $\Rp(s)>\tfrac12+\omega$ if and only if $\Psi_\omega(t)\ge 0$ for all
sufficiently large $t$. The shifted family appears to have attracted no further attention.

We make the shifted criterion quantitative. Writing $\vartheta:=\Theta-\tfrac12$ where
$\Theta=\sup\{\Rp(\rho)\}$, we show that for $0<\omega<\tfrac12$
\[
  \Psi_\omega(t)\;=\;\lxi{\bigl(\tfrac12+\omega\bigr)}\,t \;+\; \lxi{'\bigl(\tfrac12+\omega\bigr)}
  \;+\;O\!\left(e^{-(\omega-\vartheta)t}\right),
\]
where the two leading coefficients are unconditional closed forms and the implied constant is
explicit --- under the Riemann Hypothesis it is $\omega^{-1}(\xi'/\xi)(\tfrac12+\omega)$. Four
consequences follow. First, the criterion carries a \emph{margin} which is linear in $t$ with slope
exactly $(\xi'/\xi)(\tfrac12+\omega)$; this slope vanishes as $\omega\to0^{+}$ precisely because
$(\xi'/\xi)(\tfrac12)=0$, which identifies the mechanism by which the Riemann Hypothesis is the
marginal member of the family. Second, whenever the zero-free half-plane holds at all --- the only
case in which the criterion can hold --- it holds with room to spare: $\Psi_\omega(t)>0$ for every
$t>t_0(\omega)$ with $t_0(\omega)\le 0.0129\,\omega<0.0065$, beyond which $\Psi_\omega$ admits an
explicit linear lower bound. Third, a zero with $\Rp(\rho)=\tfrac12+\delta$ and $\delta>\omega$
drives $\Psi_\omega$ with an envelope growing at the exact rate $\delta-\omega$ and explicit
amplitude, so that failure of the criterion is exponentially visible. Fourth, the family undergoes
a sharp phase transition at $\omega_c:=\Theta-\tfrac12$: below $\omega_c$ the function
$e^{-\varepsilon t}\Psi_\omega(t)$ oscillates between $\pm\infty$ for every
$\varepsilon<\omega_c-\omega$, while at and above $\omega_c$ one has
$\Psi_\omega(t)=(\xi'/\xi)(\tfrac12+\omega)t+O(1)$; the Riemann Hypothesis is precisely the
statement that this transition is continuous. Unconditionally, the verification of the Riemann
Hypothesis to height $3\cdot10^{12}$, combined with the explicit zero-free regions, forces
$\Psi_\omega(t)>0$ on an explicit initial window --- for $\omega=\tfrac14$ up to $t=110$, that is,
$x=e^{t}\approx10^{48}$ --- whose endpoint obeys $t_+(\omega)=(4/c+o(1))(\tfrac12-\omega)^{-2}$
with $c=5.558691$ as $\omega\to\tfrac12^{-}$; in particular no computation of feasible size can
exhibit a violation. All results transfer to real primitive Dirichlet characters, where the margin
grows with the conductor and the criterion is provably blind to real (Siegel) zeros. We verify
each statement numerically, including checks of the asymptotics against the arithmetic sides of
the formulas that use no information about the zeros, and close with an assessment of what the
reformulation does and does not supply.
\end{abstract}

\maketitle

\setcounter{tocdepth}{1}
\tableofcontents

% ===========================================================================
\section{Introduction}
% ===========================================================================

This introduction is longer and more discursive than is customary. The results below are elementary
consequences of a formula of Suzuki, and their interest lies less in their proofs than in what they
say about the shape of a certain family of criteria for the Riemann Hypothesis. That is hard to
convey without first laying out the context, so we do.

\subsection{Zeros, primes, and the explicit formula}

The Riemann zeta-function $\zeta(s)$ and the prime numbers are linked by a duality made precise by
Riemann and von Mangoldt. In its modern form, due to Guinand \cite{Guinand} and Weil \cite{Weil},
the link is an identity: for suitable test functions $\varphi$, with
$\widehat\varphi(z):=\int_\R\varphi(x)e^{izx}\,dx$,
\begin{equation}\label{eq:explicit}
  \sum_{\gamma}\widehat{\varphi}(\gamma)
  \;=\;
  \widehat{\varphi}\!\left(\tfrac{i}{2}\right)+\widehat{\varphi}\!\left(-\tfrac{i}{2}\right)
  -\sum_{n\ge2}\frac{\Lambda(n)}{\sqrt n}\bigl(\varphi(\log n)+\varphi(-\log n)\bigr)
  - (\log\pi)\varphi(0)
  + \frac{1}{2\pi}\int_{\R}\Rp\left[\frac{\Gamma'}{\Gamma}\!\left(\tfrac14+\tfrac{iz}{2}\right)\right]
    \widehat\varphi(z)\,dz .
\end{equation}
Here $\Lambda$ is the von Mangoldt function and the sum on the left is over the zeros $\gamma$ of
$\xi(\tfrac12-iz)$, counted with multiplicity, where
\begin{equation}\label{eq:xi}
  \xi(s)=\tfrac12 s(s-1)\pi^{-s/2}\Gamma\!\left(\tfrac s2\right)\zeta(s)
\end{equation}
is the completed zeta-function, entire of order one and satisfying $\xi(s)=\xi(1-s)$. The Riemann
Hypothesis (RH) is the assertion that every zero $\rho$ of $\xi$ has $\Rp(\rho)=\tfrac12$;
equivalently, in the variable of \eqref{eq:explicit}, that every $\gamma$ is real.

Identity \eqref{eq:explicit} is an equality between a sum over zeros and a sum over prime powers.
Read from left to right it converts spectral data into arithmetic; read from right to left it does
the opposite. It is the fundamental tool for extracting information about $\zeta$ from the primes,
and conversely.

\subsection{Weil positivity, and the awkwardness of distributions}

Weil observed that \eqref{eq:explicit} encodes RH as a positivity statement. Define the
\emph{Weil distribution} $W$ by $W(\varphi):=\sum_\gamma\widehat\varphi(\gamma)$. If every $\gamma$
is real then, for $\varphi=\psi*\widetilde\psi$ with $\widetilde\psi(x)=\overline{\psi(-x)}$, one has
$\widehat\varphi(\gamma)=|\widehat\psi(\gamma)|^{2}\ge0$, so $W(\psi*\widetilde\psi)\ge0$; Weil
proved the converse. Thus
\begin{equation}\label{eq:weil}
  \text{RH}\quad\Longleftrightarrow\quad
  W(\psi*\widetilde\psi)\ge0\ \ \text{for every }\psi\in C_c^\infty(\R).
\end{equation}
This is attractive because it is \emph{local in the primes}: if $\psi$ is supported in $[-a,a]$,
only prime powers $p^m\le e^{a}$ enter the right-hand side of \eqref{eq:explicit}. An infinite
problem about all primes becomes a family of finite problems.

It is also awkward, for a reason that is easy to state: $W$ is a distribution, not a function. Its
``kernel'' is $\sum_\gamma e^{i\gamma t}$, which converges nowhere. The analytic work of Yoshida
\cite{Yoshida} and Bombieri \cite{BombieriI,BombieriII} on \eqref{eq:weil} is correspondingly
technical, and the objects that arise --- hermitian forms on spaces of smooth compactly supported
functions with side conditions --- are not the objects for which classical analysis has the most
tools.

There is a second structural difficulty. One would like to test \eqref{eq:weil} on a small family of
$\psi$ rather than on all of $C_c^\infty$, and it is a priori unclear that any small family
suffices. That it can is the content of Li's criterion \cite{Li}, which Bombieri and Lagarias
\cite{BombieriLagarias} showed to be exactly the positivity of $W$ on a distinguished countable
family of test functions. This is worth emphasising, since it is the precedent for everything below:
\emph{RH is already equivalent to Weil positivity on a one-parameter family.}

\subsection{Integrating twice: the screw function}

Suzuki's contribution \cite{Suzuki2023} is to integrate the Weil distribution twice. Set
\begin{equation}\label{eq:Psi}
  \Psi(t):=4\bigl(e^{t/2}+e^{-t/2}-2\bigr)
  -\sum_{n\le e^{t}}\frac{\Lambda(n)}{\sqrt n}(t-\log n)
  +\frac t2\left[\frac{\Gamma'}{\Gamma}\!\left(\tfrac14\right)-\log\pi\right]
  +\frac14\Bigl(C-e^{-t/2}\Phi\bigl(e^{-2t},2,\tfrac14\bigr)\Bigr)
\end{equation}
for $t\ge0$, extended to an even function of $t\in\R$; here $\Phi$ is the Hurwitz--Lerch
zeta-function and $C=\zeta(2,\tfrac14)$. Then $\Psi''=W$ as distributions, and $\Psi$ is a
\emph{continuous} real function with the two faces
\begin{equation}\label{eq:PsiTwoFaces}
  \Psi(t)=\sum_\gamma\frac{1-\cos(\gamma t)}{\gamma^{2}},
  \qquad\qquad
  \int_0^\infty \Psi(t)e^{izt}\,dt=-\frac{1}{z^{2}}\,\lxi{\left(\tfrac12-iz\right)}\quad(\Ip z>\tfrac12).
\end{equation}

Two observations make $\Psi$ concrete. First, by the prime number theorem the prime sum in
\eqref{eq:Psi} equals $4e^{t/2}+O(e^{t/2}e^{-c\sqrt t})$ \cite[Prop.~2.1]{Suzuki2023}, so the
leading exponential cancels against the first term: \emph{$\Psi$ is the error term in the prime
number theorem, twice integrated and normalised}. Second, if RH holds then every term of the series
in \eqref{eq:PsiTwoFaces} is non-negative and $\Psi$ is confined to a narrow band,
$0\le\Psi(t)<2\sum_\gamma\gamma^{-2}<0.0925$.

The resulting criterion \cite[Thm.~1.7]{Suzuki2023} is as concrete as any known form of RH:
\begin{equation}\label{eq:pointwise}
  \text{RH}\quad\Longleftrightarrow\quad \Psi(t)\ge0\ \text{ for every }t .
\end{equation}
Necessity is immediate from \eqref{eq:PsiTwoFaces}. Sufficiency does not use Weil's argument at all:
it follows from the classical theorem that the Laplace transform of a non-negative function has a
singularity at the real point of its abscissa of convergence \cite[Ch.~II, Thm.~5b]{Widder}, applied
to the second identity in \eqref{eq:PsiTwoFaces}, together with the elementary fact that
$\zeta$ has no zeros on $[\tfrac12,\infty)$.

Since $\Psi(t)=W(\Delta_t)$ where $\Delta_t=R_t*\widetilde{R_t}$ is the triangular function
supported on $[-t,t]$ and $R_t=2^{-1/2}\mathbf{1}_{[-t/2,t/2]}$ a normalised rectangular function
\cite[\S3.4]{Suzuki2023}, \eqref{eq:pointwise} says that Weil positivity tested on the family of
\emph{rectangular} functions alone already characterises RH. As noted above, the existence of such a
reduction is not new --- Li's criterion is one --- but the family here is a continuum of compactly
supported functions rather than a countable family, and the resulting object $\Psi$ is a continuous
function of one real variable with an explicit arithmetic formula. That is what makes it possible to
attach the classical theory of Kre\u\i n \cite{KreinLanger}, in which $-\Psi$ is a \emph{screw
function}: a continuous hermitian $g$ with $g(t-u)-g(t)-g(-u)+g(0)$ a non-negative kernel.
Equivalently, $-g$ is conditionally negative definite, so the natural home of the positivity is the
increments of $\Psi$, not $\Psi$ itself.

\subsection{Why RH is the hard member of a family}

Criterion \eqref{eq:pointwise} is exact, and exactly saturated. Suzuki observes
\cite[\S7.1]{Suzuki2023} that if RH holds, the zeros are simple, and their positive ordinates are
linearly independent over $\mathbb{Q}$, then Kronecker's theorem gives
\begin{equation}\label{eq:liminf}
  \limsup_{t\to\infty}\Psi(t)=2\sum_\gamma\frac{1}{\gamma^{2}},
  \qquad
  \liminf_{t\to\infty}\Psi(t)=0 .
\end{equation}
The function comes arbitrarily close to violating the inequality, infinitely often. As Suzuki puts
it, this ``suggests the difficulty of proving the non-negativity of $\Psi(t)$ by approximation.''

There is, however, a graded version. Following Lagarias \cite{Lagarias1999}, whose Theorem 1.1 is
already stated with a general abscissa $\theta$, Suzuki introduces \cite[\S11]{Suzuki2023} for
$\omega\in\R$ the shifted family
\begin{equation}\label{eq:Psiomegadef}
  \Psi_\omega(t):=e^{-\omega t}\Psi(t)+2\omega\int_0^t e^{-\omega u}\Psi(u)\,du
  +\omega^{2}\int_0^t (t-u)e^{-\omega u}\Psi(u)\,du \qquad(t>0),
\end{equation}
again extended evenly, so that $\Psi_0=\Psi$ and
\begin{equation}\label{eq:PsiomegaLaplace}
  \int_0^\infty\Psi_\omega(t)e^{izt}\,dt=-\frac1{z^{2}}\lxi{\left(\tfrac12+\omega-iz\right)},
  \qquad \Ip z>\tfrac12-\omega,
\end{equation}
and proves:

\begin{theorem}[{Suzuki \cite[Thm.~11.1]{Suzuki2023}}]\label{thm:suzuki}
Let $\omega\in\R$. Then $\xi(s)\neq0$ for $\Rp(s)>\tfrac12+\omega$ if and only if there exists
$t_0>0$ such that $\Psi_\omega(t)\ge0$ for all $t\ge t_0$.
\end{theorem}

Two remarks of Suzuki's about Theorem \ref{thm:suzuki} matter for what follows. First, in the
necessity direction the qualifier ``sufficiently large'' can be dispensed with: Suzuki observes
\cite[\S11]{Suzuki2023} that if $\xi(s)\neq0$ for $\Rp(s)>\tfrac12+\omega$ then
$i(\xi'/\xi)(\tfrac12+\omega-iz)$ is a Herglotz function, so $-\Psi_\omega$ is a screw function by
the Kre\u\i n--Langer correspondence, and $\Psi_\omega(t)\ge0$ for \emph{every} $t$. The weakening
to ``eventually'' strengthens only the sufficiency direction, where it weakens the hypothesis.
Second, the family is monotone: $\Psi_{\omega+\eta}$ is obtained from $\Psi_\omega$ by the same
transformation \eqref{eq:Psiomegadef} with parameter $\eta$, so non-negativity propagates upward in
$\omega$.

Writing $\Theta:=\sup\{\Rp(\rho):\xi(\rho)=0\}$ and
$\omega_c:=\inf\{\omega:\Psi_\omega\ge0\text{ eventually}\}$, Theorem \ref{thm:suzuki} says
$\omega_c=\Theta-\tfrac12$, and RH is the case $\omega_c=0$. Two things should be said about this
range. For $\omega\ge\tfrac12$ the criterion is vacuous, since $\zeta(s)\neq0$ for $\Rp(s)\ge1$ is
equivalent to the prime number theorem. For every $\omega<\tfrac12$ it is open: all known zero-free
regions have the shape $\sigma\ge1-f(|t|)$ with $f(|t|)\to0$ --- the Korobov--Vinogradov form
\cite{Ford}, or in explicit form $\sigma\ge1-\bigl(55.241(\log|t|)^{2/3}(\log\log|t|)^{1/3}\bigr)^{-1}$
\cite{MTY} --- and none of them contains a half-plane $\sigma>c$ with $c<1$. Whether $\Theta<1$ is
not known.

So Theorem \ref{thm:suzuki} interpolates, between a vacuous endpoint and RH, through a range of
statements every one of which is open. That makes it worth asking how the criterion behaves as a
function of $\omega$, and in particular whether the saturation \eqref{eq:liminf} that obstructs the
case $\omega=0$ persists for $\omega>0$. That is the question this note answers.

\subsection{Results}

Throughout, $\vartheta:=\Theta-\tfrac12\in[0,\tfrac12]$, and we write

\begin{hypothesis}\label{hyp}
$\xi(s)\neq0$ for $\Rp(s)>\tfrac12+\vartheta$.
\end{hypothesis}

\noindent
This holds by the definition of $\vartheta$; RH is the case $\vartheta=0$. Define
\begin{equation}\label{eq:AB}
  A(\omega):=\sum_\gamma\frac{1}{\gamma^{2}+\omega^{2}},
  \qquad
  B(\omega):=\sum_\gamma\frac{\gamma^{2}-\omega^{2}}{(\gamma^{2}+\omega^{2})^{2}},
  \qquad
  A^{*}(\omega):=\frac12\sum_\gamma\left[\frac1{|\gamma-i\omega|^{2}}+\frac1{|\gamma+i\omega|^{2}}\right].
\end{equation}

Our first observation is that the first two sums are classical objects in disguise, with no
hypothesis at all.

\begin{proposition}\label{prop:AB}
For $0<\omega<\tfrac12$ the series \eqref{eq:AB} converge absolutely, and, unconditionally,
\[
  A(\omega)=\frac1\omega\,\lxi{\left(\tfrac12+\omega\right)},
  \qquad
  B(\omega)=\lxi{'\left(\tfrac12+\omega\right)} .
\]
Moreover $A(\omega)$ is strictly decreasing on $(0,\tfrac12)$ with
$\lim_{\omega\to\frac12^{-}}A(\omega)=\sum_\rho\frac{1}{\rho(1-\rho)}=2+\gamma_0-\log4\pi
=0.0461914\ldots$, where $\gamma_0$ is Euler's constant. If RH holds, $A^{*}(\omega)=A(\omega)$.
\end{proposition}

The main result is then a two-term asymptotic with explicit error. Only the error term is
conditional.

\begin{theorem}\label{thm:asym}
Let $0<\omega<\tfrac12$ and assume $\vartheta\le\omega$. Then for all $t\ge0$,
\[
  \Psi_\omega(t)=\lxi{\left(\tfrac12+\omega\right)}\,t
  \;+\;\lxi{'\left(\tfrac12+\omega\right)}
  \;+\;R_\omega(t),
  \qquad
  |R_\omega(t)|\le A^{*}(\omega)\,e^{-(\omega-\vartheta)t}.
\]
In particular, under RH, $|R_\omega(t)|\le \omega^{-1}(\xi'/\xi)(\tfrac12+\omega)\,e^{-\omega t}$.
\end{theorem}

The linear term is what we call the \emph{margin} of the criterion, and its slope is
$(\xi'/\xi)(\tfrac12+\omega)$. Since $\xi(s)=\xi(1-s)$ forces $(\xi'/\xi)(\tfrac12)=0$, the margin
vanishes as $\omega\to0^{+}$, at the rate $(\xi'/\xi)'(\tfrac12)\,\omega=0.0462099862\ldots\times\omega$.
This identifies the mechanism behind \eqref{eq:liminf}: \emph{the marginality of the Riemann
Hypothesis inside this family is exactly the vanishing of $\xi'/\xi$ at the central point.}

The margin also shows that whenever the criterion of Theorem \ref{thm:suzuki} holds at all, it holds
with an explicit strict margin from an explicit, and very small, threshold onward.

\begin{theorem}\label{thm:threshold}
Let $0<\omega<\tfrac12$ and assume $\vartheta\le\omega$, i.e.\ $\xi(s)\neq0$ for
$\Rp(s)>\tfrac12+\omega$ --- precisely the situation in which the criterion holds. Then for every
$t\ge0$,
\[
  \Psi_\omega(t)\;\ge\;\lxi{\left(\tfrac12+\omega\right)}t+\lxi{'\left(\tfrac12+\omega\right)}
  -A^{*}(\omega),
\]
and $\Psi_\omega(t)>0$ for every $t>t_0(\omega):=\bigl(A^{*}(\omega)-B(\omega)\bigr)/\bigl(\omega
A(\omega)\bigr)$. Moreover
\[
  A^{*}(\omega)-B(\omega)\;\le\;8\,\omega^{2}\sum_\rho\frac{1}{(\Ip\rho)^{4}},
  \qquad\text{whence}\qquad
  t_0(\omega)\;\le\;\frac{8\sum_\rho(\Ip\rho)^{-4}}{2+\gamma_0-\log4\pi}\;\omega
  \;=\;0.01288\ldots\times\omega\;<\;0.0065 .
\]
\end{theorem}

The two regimes --- linear growth at and above $\omega_c:=\Theta-\tfrac12$, failure below --- meet
in a sharp transition, and the behaviour below $\omega_c$ is not merely a failure of
non-negativity but two-sided exponential oscillation.

\begin{theorem}\label{thm:transition}
Let $0<\omega<\tfrac12$.
\begin{enumerate}
\item[(i)] If $\vartheta\le\omega$ then
$\Psi_\omega(t)=(\xi'/\xi)(\tfrac12+\omega)\,t+O(1)$ with $|O(1)|\le|B(\omega)|+A^{*}(\omega)$;
in particular $\Psi_\omega(t)/t\to(\xi'/\xi)(\tfrac12+\omega)$.
\item[(ii)] If $\omega<\vartheta$ then for every $0\le\varepsilon<\vartheta-\omega$,
\[
  \limsup_{t\to\infty}e^{-\varepsilon t}\Psi_\omega(t)=+\infty,
  \qquad
  \liminf_{t\to\infty}e^{-\varepsilon t}\Psi_\omega(t)=-\infty .
\]
\end{enumerate}
\end{theorem}

\begin{corollary}\label{cor:graded}
$\Theta\le\tfrac12+\omega$ if and only if
$\Psi_\omega(t)=(\xi'/\xi)(\tfrac12+\omega)\,t+O(1)$. Consequently
$m(\omega):=\liminf_{t\to\infty}\Psi_\omega(t)/t$ equals $-\infty$ on $(0,\omega_c)$ and
$(\xi'/\xi)(\tfrac12+\omega)$ on $[\omega_c,\tfrac12)$: the normalised slope is an order
parameter, with a jump of size at least $(2+\gamma_0-\log4\pi)\,\omega_c$ at $\omega_c$. The
Riemann Hypothesis holds if and only if the transition is continuous, i.e.\ $\omega_c=0$ and
$m(\omega)\to0$ as $\omega\to0^{+}$.
\end{corollary}

Corollary \ref{cor:graded} is a graded form of Suzuki's Theorem 1.6 (RH $\Leftrightarrow$
$\Psi=O(1)$), with the sharp linear term; it also yields yet another description of the critical
shift: $\omega_c=\inf\{\omega:\Psi_\omega(t)-(\xi'/\xi)(\tfrac12+\omega)t\text{ is bounded}\}$.

Next, failure of the criterion is exponentially visible, at a rate that reads off the depth of
the violation.

\begin{proposition}\label{prop:sens}
Let $\rho_0=\tfrac12+\delta+i\tau_0$ be a zero of $\xi$ with $\delta>\omega\ge0$ and $\tau_0>0$, and
let $\Psi_\omega^{(\rho_0)}$ denote the contribution to the series for $\Psi_\omega$ of the quartet
$\{\rho_0,\overline{\rho_0},1-\rho_0,1-\overline{\rho_0}\}$. Then
\[
  \Psi_\omega^{(\rho_0)}(t)
  = -\,\frac{2\,e^{(\delta-\omega)t}}{\tau_0^{2}+(\delta-\omega)^{2}}\,
     \cos\bigl(\tau_0 t+\alpha\bigr)\;+\;O(t),
  \qquad \alpha=2\arg\bigl(\tau_0+i(\delta-\omega)\bigr);
\]
in particular
$\limsup_{t\to\infty}e^{-(\delta-\omega)t}\bigl|\Psi_\omega^{(\rho_0)}(t)\bigr|
 =2\bigl(\tau_0^{2}+(\delta-\omega)^{2}\bigr)^{-1}>0$.
\end{proposition}

Combining the sensitivity estimate with the numerical verification of RH to height $3\cdot10^{12}$
\cite{PlattTrudgian} gives an unconditional initial window of positivity, which we record because it
delimits sharply what computation can and cannot say about the criterion.

\begin{proposition}\label{prop:window}
Let $0<\omega<\tfrac12$ and $T_v=3\cdot10^{12}$. Unconditionally, for every $t\ge0$,
\[
  \Psi_\omega(t)\;\ge\;\lxi{\left(\tfrac12+\omega\right)}t+\lxi{'\left(\tfrac12+\omega\right)}
  -A(\omega)\,e^{-\omega t}-e^{(\frac12-\omega)t}\,E(t),
\]
where, with $L_v:=\log T_v$ and
$\nu(L):=\max\bigl\{(5.558691\,L)^{-1},\,(55.241\,L^{2/3}(\log L)^{1/3})^{-1}\bigr\}$ the explicit
zero-free function of \cite{MTY},
\[
  E(t):=\frac1\pi\int_{L_v}^{\infty}(L-\log2\pi)\,e^{-L-\nu(L)\,t}\,dL,
  \qquad
  E(t)\le E(0)=S:=\!\!\sum_{|\Ip\rho|>T_v}\!\!\frac{1}{(\Ip\rho)^{2}}\;<\;3.2\cdot10^{-12}.
\]
Consequently $\Psi_\omega(t)>0$ on an explicit window $t_{-}(\omega)<t<t_{+}(\omega)$ with
$t_{-}(\omega)<0.0017\,\omega$ and, for example, $t_{+}(0.1)=64.7$, $t_{+}(0.25)=109.9$,
$t_{+}(0.4)=300.9$, $t_{+}(0.49)=7262$; as $\omega\to\tfrac12^{-}$,
\[
  t_{+}(\omega)=\Bigl(\frac4c+o(1)\Bigr)\Bigl(\tfrac12-\omega\Bigr)^{-2},
  \qquad c=5.558691 .
\]
\end{proposition}

All of this transfers to real primitive Dirichlet characters, with two instructive changes
(Section \ref{sec:dirichlet}). The closed forms, the margin, the threshold and the transition hold
verbatim for the shifted screw functions $\Psi_{\omega,\chi}$, with slope
$(\xi'/\xi)(\tfrac12+\omega,\chi)$ vanishing at $\omega=0$ for the same functional-equation reason
--- and the margin \emph{grows} with the conductor ($B_\chi(0)=0.113,\ 0.156,\ 0.157$ for
$q=3,4,5$, against $0.0462$ for $\zeta$): within this family, $\zeta$ is the most marginal member.
The arithmetic side is Proposition \ref{prop:chiprime}. The one genuinely new phenomenon is
negative: the criterion is blind to real zeros --- a Siegel zero pushes $\Psi_{\omega,\chi}$
\emph{up}, exponentially (Proposition \ref{prop:siegel}) --- so what the shifted criterion
characterises for $\chi$ is the location of the \emph{non-real} zeros only. This explains
quantitatively the real-zero hypothesis in Suzuki's extension of the screw-function theory to the
extended Selberg class \cite[Thm.~1.1]{SuzukiSelberg}.

Sections \ref{sec:setup}--\ref{sec:sens} prove these statements; Section \ref{sec:dirichlet}
carries them to Dirichlet $L$-functions and discusses the extended Selberg class; Section
\ref{sec:numerics} verifies them numerically, including a test of Theorem \ref{thm:asym} against
the arithmetic side of \eqref{eq:Psi} which uses no information about the zeros; Section
\ref{sec:discussion} assesses what the reformulation supplies.

\subsection*{Attribution}
The function $\Psi$, the shifted family $\Psi_\omega$, Theorem \ref{thm:suzuki}, and the observation
that a zero-free half-plane forces $\Psi_\omega\ge0$ pointwise are Suzuki's \cite{Suzuki2023}; the
screw-function theory over the extended Selberg class, including the role of real zeros, is
developed by him in \cite{SuzukiSelberg}. The graded criterion in the form ``$\Rp(\rho)\le\theta$
for all $\rho$ if and only if a certain function is Herglotz on $\Rp(s)>\theta$'' is Lagarias's
\cite[Thm.~1.1]{Lagarias1999}. The reduction of RH to Weil positivity on a one-parameter family of
test functions goes back to Li \cite{Li} and Bombieri--Lagarias \cite{BombieriLagarias}. What is
offered here is Proposition \ref{prop:AB}, Theorems \ref{thm:asym}, \ref{thm:threshold} and
\ref{thm:transition}, Corollary \ref{cor:graded}, Propositions \ref{prop:sens} and
\ref{prop:window}, the Dirichlet and Selberg-class material of Section \ref{sec:dirichlet}
(Theorem \ref{thm:chi}, Propositions \ref{prop:chiprime} and \ref{prop:siegel}), the numerical
verification, and the discussion in Section \ref{sec:discussion}.

% ===========================================================================
\section{Setup}\label{sec:setup}
% ===========================================================================

We use Suzuki's normalisation throughout. Let $\Xi(z):=\xi(\tfrac12-iz)$, an entire function of
order one; by $\xi(s)=\xi(1-s)$ it is even, and its zeros $\gamma$ are related to the zeros $\rho$
of $\xi$ by
\begin{equation}\label{eq:gammarho}
  \gamma=i\left(\rho-\tfrac12\right),
  \qquad\text{i.e.}\qquad
  \rho=\beta+i\tau \;\Longrightarrow\; \gamma=-\tau+i\bigl(\beta-\tfrac12\bigr).
\end{equation}
Thus $\Ip\gamma=\Rp\rho-\tfrac12$, RH is the statement that all $\gamma$ are real, and
Hypothesis \ref{hyp} says $|\Ip\gamma|\le\vartheta$ for every $\gamma$. Zeros come in the symmetric
configuration $\gamma\mapsto-\gamma$, $\gamma\mapsto\overline{\gamma}$. Two unconditional facts
about the zeros are used repeatedly: $|\Ip\gamma|\le\tfrac12$ for every $\gamma$
\cite[Thm.~2.12]{Titchmarsh}, and $|\Rp\gamma|\ge14$, since the nontrivial zero of least height has
ordinate $14.1347\ldots$ \cite[\S15]{Titchmarsh}.

We shall use the standard Hadamard consequences (see e.g.\ \cite[\S12]{Davenport})
\begin{equation}\label{eq:hadamard}
  \frac{\Xi'}{\Xi}(z)=\sum_\gamma\frac{1}{z-\gamma},
  \qquad
  \frac{\Xi'}{\Xi}(z)=-i\,\lxi{\left(\tfrac12-iz\right)},
  \qquad
  \lxi{(1-s)}=-\lxi{(s)},
\end{equation}
the first with the usual convention that $\gamma$ and $-\gamma$ are paired.

The starting point is Suzuki's series for the shifted family \cite[\S11]{Suzuki2023}: for $t>0$,
\begin{equation}\label{eq:Psiomegaseries}
  \Psi_\omega(t)=\sum_\gamma\frac{1}{(\gamma^{2}+\omega^{2})^{2}}
  \Bigl[\;\omega t(\gamma^{2}+\omega^{2})+(\gamma^{2}-\omega^{2})
   - e^{-\omega t}(\gamma^{2}-\omega^{2})\cos(\gamma t)
   - e^{-\omega t}\,2\gamma\omega\sin(\gamma t)\;\Bigr].
\end{equation}
Setting $\omega=0$ recovers $\Psi(t)=\sum_\gamma\gamma^{-2}(1-\cos\gamma t)$, as it must. Suzuki
also gives the arithmetic counterpart, which we record because it is what we compute with in
Section \ref{sec:numerics}: with $\psi=\Gamma'/\Gamma$ and $\psi^{(1)}=\psi'$,
\begin{multline}\label{eq:Psiomegaprimes}
  \Psi_\omega(t)=
  4\left[\frac{e^{(\frac12-\omega)t}}{(1-2\omega)^{2}}+\frac{e^{-(\frac12+\omega)t}}{(1+2\omega)^{2}}
  -\frac{4-2(1-\omega t)(1-4\omega^{2})}{(1-4\omega^{2})^{2}}\right]
  +\frac t2\left[\psi\!\left(\tfrac14+\tfrac\omega2\right)-\log\pi\right]\\
  +\frac14\left[\psi^{(1)}\!\left(\tfrac14+\tfrac\omega2\right)
   -e^{-(\frac12+\omega)t}\,\Phi\!\left(e^{-2t},2,\tfrac14+\tfrac\omega2\right)\right]
  -\sum_{n\le e^{t}}\frac{\Lambda(n)}{n^{\frac12+\omega}}(t-\log n).
\end{multline}

\begin{lemma}\label{lem:conv}
Let $0<\omega<\tfrac12$. Unconditionally,
$\Rp(\gamma^{2}+\omega^{2})\ge(\Rp\gamma)^{2}-\tfrac14+\omega^{2}>0$ for every $\gamma$, the series
\eqref{eq:AB} converge absolutely, and
$\sum_\gamma|\gamma|\,|\gamma^{2}+\omega^{2}|^{-2}<\infty$.
\end{lemma}

\begin{proof}
Write $\gamma=a+ib$; then $|b|\le\tfrac12$ and $|a|\ge14$. We have
$\gamma^{2}+\omega^{2}=a^{2}-b^{2}+\omega^{2}+2iab$, whose real part is at least
$a^{2}-\tfrac14>0$. Hence $|\gamma^{2}+\omega^{2}|\ge a^{2}-\tfrac14$, and likewise
$|\gamma\pm i\omega|^{2}=a^{2}+(b\pm\omega)^{2}\ge a^{2}$. By the Riemann--von Mangoldt formula the
number of $\gamma$ with $|a|\le T$ is $O(T\log T)$, so $\sum_\gamma(a^{2}-\tfrac14)^{-1}<\infty$ and
$\sum_\gamma|\gamma|(a^{2}-\tfrac14)^{-2}<\infty$ by partial summation. The stated bounds follow.
\end{proof}

\begin{remark}
Suzuki also records the composition law: $\Psi_{\omega+\eta}$ arises from $\Psi_\omega$ by the
transformation \eqref{eq:Psiomegadef} with parameter $\eta>0$ \cite[\S11]{Suzuki2023}. Since that
transformation has non-negative kernel, non-negativity of $\Psi_\omega$ (pointwise, or on $[0,t_1)$)
implies the same for $\Psi_{\omega'}$, $\omega'>\omega$: the family is monotone in $\omega$.
\end{remark}

% ===========================================================================
\section{Two spectral sums in closed form}\label{sec:AB}
% ===========================================================================

\begin{proof}[Proof of Proposition \ref{prop:AB}]
Convergence is Lemma \ref{lem:conv}, and requires no hypothesis. For the closed forms, evaluate the
first identity of \eqref{eq:hadamard} at $z=\pm i\omega$ (legitimate for $0<\omega<\tfrac12$ since
no $\gamma$ lies on the imaginary axis: $|\Rp\gamma|\ge14$) and use the second:
\[
  \sum_\gamma\frac{1}{i\omega-\gamma}=-i\,\lxi{\left(\tfrac12+\omega\right)},
  \qquad
  \sum_\gamma\frac{1}{-i\omega-\gamma}=-i\,\lxi{\left(\tfrac12-\omega\right)}
  \;=\;i\,\lxi{\left(\tfrac12+\omega\right)},
\]
the last step by the functional equation in \eqref{eq:hadamard}. Therefore
\[
  \sum_\gamma\frac{1}{\gamma-i\omega}=i\,\lxi{\left(\tfrac12+\omega\right)},
  \qquad
  \sum_\gamma\frac{1}{\gamma+i\omega}=-i\,\lxi{\left(\tfrac12+\omega\right)} .
\]
Since $\gamma^{2}+\omega^{2}=(\gamma-i\omega)(\gamma+i\omega)$, partial fractions give
\[
  A(\omega)=\sum_\gamma\frac{1}{(\gamma-i\omega)(\gamma+i\omega)}
  =\frac{1}{2i\omega}\sum_\gamma\left[\frac{1}{\gamma-i\omega}-\frac{1}{\gamma+i\omega}\right]
  =\frac{1}{2i\omega}\cdot 2i\,\lxi{\left(\tfrac12+\omega\right)}
  =\frac1\omega\lxi{\left(\tfrac12+\omega\right)},
\]
the rearrangement being justified by absolute convergence of $A(\omega)$ together with the standard
pairing of $\gamma$ with $-\gamma$ in \eqref{eq:hadamard}.

For $B$, note the elementary identity
\[
  \frac{\partial}{\partial\omega}\left[\frac{\omega}{\gamma^{2}+\omega^{2}}\right]
  =\frac{\gamma^{2}-\omega^{2}}{(\gamma^{2}+\omega^{2})^{2}} .
\]
By Lemma \ref{lem:conv} the differentiated series converges locally uniformly on
$0<\omega<\tfrac12$, so term-by-term differentiation of $\omega A(\omega)$ is legitimate and
\[
  B(\omega)=\frac{d}{d\omega}\bigl[\omega A(\omega)\bigr]
  =\frac{d}{d\omega}\lxi{\left(\tfrac12+\omega\right)}
  =\lxi{'\left(\tfrac12+\omega\right)} .
\]

For monotonicity, pair each $\gamma$ with $\overline\gamma$, so that
$A(\omega)=\sum_\gamma X_\gamma/(X_\gamma^{2}+Y_\gamma^{2})$ with
$X_\gamma=a^{2}-b^{2}+\omega^{2}>0$ and $Y_\gamma=2ab$. Then
\[
  \frac{\partial}{\partial(\omega^{2})}\,\frac{X}{X^{2}+Y^{2}}
  =\frac{Y^{2}-X^{2}}{(X^{2}+Y^{2})^{2}}<0,
\]
since $X>|Y|$ is equivalent to $(|a|-|b|)^{2}+\omega^{2}>0$. So every paired term, hence
$A(\omega)$, is strictly decreasing. At $\omega=\tfrac12$, $\gamma^{2}+\tfrac14=\rho(1-\rho)$ by
\eqref{eq:gammarho}, and $\sum_\rho 1/(\rho(1-\rho))=2\sum_\rho 1/\rho=2+\gamma_0-\log4\pi$ is
classical \cite[\S12]{Davenport}; the limit follows by dominated convergence. Finally, if RH holds
then $\gamma$ is real, $|\gamma\pm i\omega|^{2}=\gamma^{2}+\omega^{2}$, and $A^{*}=A$.
\end{proof}

\begin{remark}\label{rem:B0}
Letting $\omega\to0^{+}$ gives
$B(0)=\sum_\gamma\gamma^{-2}=(\xi'/\xi)'(\tfrac12)=0.0462099862308\ldots$ (an unconditional
identity), so $2B(0)=0.0924199725\ldots$, consistent with Suzuki's conditional bound
$\sup\Psi<0.094$. Also $A(\omega)\to B(0)$ as $\omega\to0^{+}$, since $(\xi'/\xi)(\tfrac12)=0$.
Note the coincidence of scales: $A$ decreases only from $0.046210$ to $0.046191$ across the whole
interval $(0,\tfrac12)$.
\end{remark}

% ===========================================================================
\section{The asymptotic}\label{sec:asym}
% ===========================================================================

\begin{proof}[Proof of Theorem \ref{thm:asym}]
Split \eqref{eq:Psiomegaseries} into the two groups of terms. The first two terms sum, by
\eqref{eq:AB} and Proposition \ref{prop:AB}, to
\[
  \omega t\sum_\gamma\frac{1}{\gamma^{2}+\omega^{2}}
  +\sum_\gamma\frac{\gamma^{2}-\omega^{2}}{(\gamma^{2}+\omega^{2})^{2}}
  =\omega A(\omega)\,t+B(\omega)
  =\lxi{\left(\tfrac12+\omega\right)}t+\lxi{'\left(\tfrac12+\omega\right)},
\]
which is the stated main term, and is unconditional. The remainder is
\begin{equation}\label{eq:R}
  R_\omega(t)=-e^{-\omega t}\sum_\gamma
  \frac{(\gamma^{2}-\omega^{2})\cos(\gamma t)+2\gamma\omega\sin(\gamma t)}
       {(\gamma^{2}+\omega^{2})^{2}} .
\end{equation}
The key is the exact decomposition
\begin{equation}\label{eq:Rdecomp}
  (\gamma^{2}-\omega^{2})\cos(\gamma t)+2\gamma\omega\sin(\gamma t)
  \;=\;\tfrac12\,e^{i\gamma t}(\gamma-i\omega)^{2}\;+\;\tfrac12\,e^{-i\gamma t}(\gamma+i\omega)^{2},
\end{equation}
which is immediate from $\cos w=\tfrac12(e^{iw}+e^{-iw})$, $\sin w=\tfrac1{2i}(e^{iw}-e^{-iw})$ and
$(\gamma\mp i\omega)^{2}=\gamma^{2}-\omega^{2}\mp2i\gamma\omega$. Dividing by
$(\gamma^{2}+\omega^{2})^{2}=(\gamma-i\omega)^{2}(\gamma+i\omega)^{2}$,
\begin{equation}\label{eq:Rclean}
  R_\omega(t)=-\frac{e^{-\omega t}}2\sum_\gamma
  \left[\frac{e^{i\gamma t}}{(\gamma+i\omega)^{2}}+\frac{e^{-i\gamma t}}{(\gamma-i\omega)^{2}}\right].
\end{equation}
Writing $\gamma=a+ib$ with $|b|\le\vartheta$, we have $|e^{\pm i\gamma t}|=e^{\mp bt}\le
e^{\vartheta t}$, whence
\[
  |R_\omega(t)|\le \frac{e^{-(\omega-\vartheta) t}}{2}\sum_\gamma
  \left[\frac{1}{|\gamma+i\omega|^{2}}+\frac{1}{|\gamma-i\omega|^{2}}\right]
  =A^{*}(\omega)\,e^{-(\omega-\vartheta)t}.
\]
Under RH, $A^{*}=A$ by Proposition \ref{prop:AB}, giving the stated constant.
\end{proof}

\begin{remark}
The one-line amplitude computation behind \eqref{eq:Rdecomp} --- for real $\gamma$,
$(\gamma^{2}-\omega^{2})\cos\theta+2\gamma\omega\sin\theta$ has amplitude exactly
$\sqrt{(\gamma^{2}-\omega^{2})^{2}+4\gamma^{2}\omega^{2}}=\gamma^{2}+\omega^{2}$ --- shows the
constant $A(\omega)$ in the RH case is the best possible for a term-by-term estimate. Note also that
the bound holds at $\vartheta=\omega$: there the remainder no longer decays, but stays bounded by
$A^{*}(\omega)$, which is all Theorem \ref{thm:threshold} needs.
\end{remark}

\begin{corollary}\label{cor:margin}
The slope of the linear term in Theorem \ref{thm:asym} is $(\xi'/\xi)(\tfrac12+\omega)$.
Unconditionally, it is positive for $0<\omega<\tfrac12$, indeed
\[
  \lxi{\left(\tfrac12+\omega\right)}=\omega A(\omega)\ \ge\ (2+\gamma_0-\log4\pi)\,\omega
  =0.0461914\ldots\times\omega,
\]
it vanishes at $\omega=0$, and it satisfies
$(\xi'/\xi)(\tfrac12+\omega)=B(0)\,\omega+O(\omega^{3})$ with $B(0)=0.04620998623\ldots$.
\end{corollary}

\begin{proof}
$\xi(s)=\xi(1-s)$ gives $(\xi'/\xi)(1-s)=-(\xi'/\xi)(s)$, so $(\xi'/\xi)$ is odd about
$s=\tfrac12$; in particular $(\xi'/\xi)(\tfrac12)=0$ and its Taylor expansion there contains only
odd powers. The linear coefficient is $(\xi'/\xi)'(\tfrac12)=B(0)$ by Remark \ref{rem:B0}. The
lower bound is the monotonicity statement of Proposition \ref{prop:AB}.
\end{proof}

\begin{remark}
Corollary \ref{cor:margin} is the precise form of the statement that RH is the marginal member of
the family. The obstruction \eqref{eq:liminf} --- $\liminf\Psi=0$, so that no approximation argument
can establish $\Psi\ge0$ --- is the $\omega=0$ shadow of a margin that is strictly positive and
grows linearly in $t$ for every $\omega>0$. Whatever the difficulty of proving
$\Psi_\omega\ge0$ for a given $\omega\in(0,\tfrac12)$, it is not that the inequality is
asymptotically saturated. Theorem \ref{thm:transition} sharpens this into a dichotomy.
\end{remark}

% ===========================================================================
\section{The threshold}\label{sec:threshold}
% ===========================================================================

\begin{proof}[Proof of Theorem \ref{thm:threshold}]
Under $\vartheta\le\omega$ we have $e^{-(\omega-\vartheta)t}\le1$ for $t\ge0$, so Theorem
\ref{thm:asym} gives the uniform lower bound
$\Psi_\omega(t)\ge\omega A(\omega)t+B(\omega)-A^{*}(\omega)$, and strict positivity for
$t>t_0=(A^{*}-B)/(\omega A)$, since $A(\omega)>0$.

It remains to bound $A^{*}-B$. Write $z_1=\gamma-i\omega$, $z_2=\gamma+i\omega$. Then
$z_1^{2}+z_2^{2}=2(\gamma^{2}-\omega^{2})$ and $z_1^{2}z_2^{2}=(\gamma^{2}+\omega^{2})^{2}$, so the
$B$-term of each zero is
\[
  \frac{\gamma^{2}-\omega^{2}}{(\gamma^{2}+\omega^{2})^{2}}
  =\frac{z_1^{2}+z_2^{2}}{2\,z_1^{2}z_2^{2}}
  =\frac12\left[\frac{1}{z_1^{2}}+\frac{1}{z_2^{2}}\right],
\]
while the $A^{*}$-term is $\tfrac12(|z_1|^{-2}+|z_2|^{-2})$. For any $z=x+iy\neq0$,
\[
  \frac{1}{|z|^{2}}-\Rp\frac{1}{z^{2}}
  =\frac{(x^{2}+y^{2})-(x^{2}-y^{2})}{|z|^{4}}
  =\frac{2y^{2}}{|z|^{4}} .
\]
Summing over the zeros (the imaginary parts cancel in conjugate pairs, so $B$ equals the sum of the
real parts) and noting $\Ip z_1=b-\omega$, $\Ip z_2=b+\omega$,
\begin{equation}\label{eq:AstarminusB}
  A^{*}(\omega)-B(\omega)
  =\sum_\gamma\left[\frac{(b-\omega)^{2}}{|\gamma-i\omega|^{4}}
  +\frac{(b+\omega)^{2}}{|\gamma+i\omega|^{4}}\right].
\end{equation}
Under $|b|\le\vartheta\le\omega$ we have $(b\mp\omega)^{2}\le4\omega^{2}$ and
$|\gamma\mp i\omega|^{2}\ge a^{2}$, whence
$A^{*}-B\le8\omega^{2}\sum_\gamma a^{-4}=8\omega^{2}\sum_\rho(\Ip\rho)^{-4}$. Computing the first
$120$ ordinates and estimating the tail with the Riemann--von Mangoldt formula (with the explicit
error bound of \cite{TrudgianArg}, whose contribution is far below the last digit retained) gives
$\sum_\rho(\Ip\rho)^{-4}=7.4346\ldots\cdot10^{-5}<7.44\cdot10^{-5}$. Together with
$\omega A(\omega)=(\xi'/\xi)(\tfrac12+\omega)\ge(2+\gamma_0-\log4\pi)\,\omega$ from Corollary
\ref{cor:margin},
\[
  t_0(\omega)\le\frac{8\cdot7.44\cdot10^{-5}}{0.0461914}\,\omega=0.01288\ldots\times\omega,
\]
which is less than $0.0065$ throughout $0<\omega<\tfrac12$.
\end{proof}

\begin{remark}\label{rem:trivial}
Two cheaper routes to non-strict positivity should be set beside Theorem \ref{thm:threshold}, both
already in \cite{Suzuki2023}. Under RH, $\Psi\ge0$ pointwise, and every term of the definition
\eqref{eq:Psiomegadef} is then non-negative, so $\Psi_\omega\ge0$ everywhere --- indeed
$\Psi_\omega(t)>0$ for $t>0$, since $\Psi(t)>0$ for $t\neq0$ \cite[Thm.~1.7]{Suzuki2023}. More
generally, under $\vartheta\le\omega$ alone, Suzuki's Herglotz argument (\S1.4 above) gives
$\Psi_\omega\ge0$ everywhere via Kre\u\i n--Langer. What Theorem \ref{thm:threshold} adds is
quantitative and elementary: an explicit linear lower bound with an explicit onset, obtained
without the Kre\u\i n--Langer correspondence. In particular the criterion, wherever it holds, holds
\emph{with room to spare} --- there is no analogue at $\omega>0$ of the saturation
\eqref{eq:liminf}.
\end{remark}

\begin{remark}\label{rem:detector}
The contrapositive is a single-point detector: a single verified value $\Psi_\omega(t^{*})\le0$ at
any $t^{*}>0.0129\,\omega$ certifies a zero with $\Rp\rho>\tfrac12+\omega$ (by Suzuki's pointwise
statement, any $t^{*}>0$ in fact suffices; the present form has the advantage of an explicit
quantitative margin below which $\Psi_\omega$ cannot dip). Proposition \ref{prop:sens} quantifies
the rate at which such values appear when a violating zero exists. Proposition \ref{prop:window}
shows, however, that no feasible computation can reach the regime where such values could occur.
\end{remark}

\begin{remark}
The threshold $t_0(\omega)\le0.0129\,\omega$ tends to $0$ with $\omega$, and covers the boundary
case $\vartheta=\omega$, where the oscillation is undamped: the reason a bounded, non-decaying
oscillation cannot break the criterion is visible in \eqref{eq:AstarminusB} --- a zero at height
$\tau$ on the boundary $\Rp\rho=\tfrac12+\omega$ contributes oscillation of amplitude
$\approx2/\tau^{2}$, but contributes $\approx4/\tau^{2}$ to the intercept $B(\omega)$; its own mean
covers its own oscillation twice over. Only zeros \emph{beyond} the half-plane can drive
$\Psi_\omega$ negative for large $t$, which is Proposition \ref{prop:sens}.
\end{remark}

% ===========================================================================
\section{The transition at $\omega_c$}\label{sec:transition}
% ===========================================================================

Write $\omega_c:=\Theta-\tfrac12=\vartheta$, so that Theorem \ref{thm:suzuki} reads: the criterion
at level $\omega$ holds if and only if $\omega\ge\omega_c$. Theorems \ref{thm:asym} and
\ref{thm:threshold} describe $\Psi_\omega$ at and above $\omega_c$; Theorem \ref{thm:transition}
adds the description below, by the same Laplace-transform mechanism that proves the sufficiency
direction of Theorem \ref{thm:suzuki}.

\begin{proof}[Proof of Theorem \ref{thm:transition}]
Part (i) is Theorem \ref{thm:asym}: for $\vartheta\le\omega$ the remainder is bounded by
$A^{*}(\omega)$.

For part (ii), suppose $\omega<\vartheta$ and, for a contradiction, that there are constants
$C,D>0$ and $0\le\varepsilon<\vartheta-\omega$ with $\Psi_\omega(t)\ge-Ce^{\varepsilon t}-D$ for
all $t\ge0$; set $f:=\Psi_\omega+Ce^{\varepsilon t}+D\ge0$. From \eqref{eq:Rclean},
\[
  |\Psi_\omega(t)|\;\le\;\omega A(\omega)\,t+|B(\omega)|+A^{*}(\omega)\,e^{(\vartheta-\omega)t},
\]
an estimate that requires no ordering of $\vartheta$ and $\omega$ (and $A^{*}$ converges
unconditionally by Lemma \ref{lem:conv}). Hence the Laplace transform
$F(z)=\int_0^\infty f(t)e^{izt}\,dt$ converges, and is analytic, for $\Ip z>\vartheta-\omega$,
and equals there
\begin{equation}\label{eq:Ftransform}
  F(z)=-\frac{1}{z^{2}}\lxi{\left(\tfrac12+\omega-iz\right)}
  +\frac{iC}{z-i\varepsilon}+\frac{iD}{z}:
\end{equation}
on $\Ip z>\tfrac12-\omega$ this is \eqref{eq:PsiomegaLaplace} together with two elementary
transforms, and the identity extends to $\Ip z>\vartheta-\omega$ by analytic continuation.
Let $\sigma_c\le\vartheta-\omega$ denote the abscissa of convergence of $F$. Since $f\ge0$,
Landau's theorem \cite[Ch.~II, Thm.~5b]{Widder} makes $i\sigma_c$ a singularity of $F$: no
analytic continuation to a neighbourhood of $i\sigma_c$ exists. Now the right-hand side of
\eqref{eq:Ftransform} is meromorphic on $\C$, and its only possible singularities on the ray
$\{iy:y>\varepsilon\}$ are poles of $(\xi'/\xi)$ at real zeros $\tfrac12+\omega+y\in(\tfrac12,1]$
--- and there are none. Hence $\sigma_c\le\varepsilon$, so $F$ is analytic on
$\Ip z>\varepsilon$ and coincides there with \eqref{eq:Ftransform} by analytic continuation;
consequently \eqref{eq:Ftransform} has no poles with $\Ip z>\varepsilon$, i.e.\ every zero
satisfies $\Rp\rho\le\tfrac12+\omega+\varepsilon$. Since $\omega+\varepsilon<\vartheta$, this
contradicts the definition of $\vartheta$. Therefore no such $C,D$ exist, which is the
$\liminf$ statement. The $\limsup$ statement follows from the same argument applied to
$f:=Ce^{\varepsilon t}+D-\Psi_\omega$, whose transform is \eqref{eq:Ftransform} with the sign of
the first term reversed.
\end{proof}

\begin{proof}[Proof of Corollary \ref{cor:graded}]
If $\Theta\le\tfrac12+\omega$, boundedness of $\Psi_\omega-(\xi'/\xi)(\tfrac12+\omega)t$ is
Theorem \ref{thm:transition}(i). Conversely, if it is bounded then
$e^{-\varepsilon t}\Psi_\omega(t)\to0$ for every $\varepsilon>0$, so (ii) fails, forcing
$\vartheta\le\omega$. For the order parameter: on $[\omega_c,\tfrac12)$, $m$ equals the slope by
(i); for $\omega<\omega_c$, part (ii) with any $0<\varepsilon<\omega_c-\omega$ yields points
$t_n\to\infty$ with $\Psi_\omega(t_n)\le-e^{\varepsilon t_n}$, so
$\Psi_\omega(t_n)/t_n\to-\infty$ and $m(\omega)=-\infty$. If $\omega_c>0$, then
$m(\omega_c)=(\xi'/\xi)(\tfrac12+\omega_c)\ge(2+\gamma_0-\log4\pi)\,\omega_c>0$ by Corollary
\ref{cor:margin}, while $m=-\infty$ immediately below: a jump. Under RH, $\omega_c=0$ and
$m(\omega)=(\xi'/\xi)(\tfrac12+\omega)\to0$ as $\omega\to0^{+}$, matching
$\liminf\Psi(t)/t=0$ at $\omega=0$ ($\Psi$ is bounded under RH).
\end{proof}

\begin{remark}
Below $\omega_c$ the oscillation is genuinely two-sided and of exponential rate arbitrarily close
to $\omega_c-\omega$; at and above $\omega_c$ the function is linear plus $O(1)$. In the language
of phase transitions, $m(\omega)$ is an order parameter and RH asserts a second-order (continuous)
transition; any $\omega_c>0$ would make it first-order, with ``latent margin''
$(\xi'/\xi)(\tfrac12+\omega_c)\approx0.0462\,\omega_c$. For $\zeta$ the proof of part (ii) uses
only the absence of real zeros in $(\tfrac12,1]$; for Dirichlet $L$-functions the same argument is
sensitive to real zeros, and Section \ref{sec:dirichlet} turns that sensitivity into a statement.
\end{remark}

% ===========================================================================
\section{An unconditional window}\label{sec:window}
% ===========================================================================

\begin{proof}[Proof of Proposition \ref{prop:window}]
Split \eqref{eq:Rclean} at the verification height: by \cite{PlattTrudgian}, every zero with
$|\Ip\rho|\le T_v=3\cdot10^{12}$ has $\Rp\rho=\tfrac12$, i.e.\ the corresponding $\gamma$ are real.
For those zeros the summand of \eqref{eq:Rclean} is bounded by $(\gamma^{2}+\omega^{2})^{-1}$, and
since the omitted high zeros have $\Rp(\gamma^{2}+\omega^{2})>0$ (Lemma \ref{lem:conv}), the partial
sum is at most $A(\omega)$; their total contribution to $R_\omega$ is at most $A(\omega)e^{-\omega
t}$ in absolute value. For a zero with $|\Ip\rho|=|a|>T_v$, the explicit zero-free regions of
\cite{MTY} give $|b|\le\tfrac12-\nu(\log|a|)$ with $\nu$ as in the statement (both regions are
valid, so their maximum is), and $|\gamma\pm i\omega|^{2}\ge a^{2}$: its contribution is at most
$e^{-\omega t}\,e^{(\frac12-\nu(\log|a|))t}\,a^{-2}$. Summing against the Riemann--von Mangoldt
density $dN(u)=(2\pi)^{-1}\log(u/2\pi)\,du$ and substituting $L=\log u$ gives
$|R_\omega(t)|\le A(\omega)e^{-\omega t}+e^{(\frac12-\omega)t}E(t)$; the explicit error bound of
\cite{TrudgianArg} for the zero-counting function perturbs $E(t)$ by less than $10^{-23}$, which
is absorbed in every stated numerical value. At $t=0$,
\[
  E(0)=S=2\int_{T_v}^{\infty}\frac{dN(u)}{u^{2}}
  =\frac{1}{\pi}\cdot\frac{\log(T_v/2\pi)+1}{T_v}+O\!\left(\frac{\log T_v}{T_v^{2}}\right)
  <3.2\cdot10^{-12},
\]
and $E$ is decreasing in $t$. The displayed lower bound for $\Psi_\omega$ follows from Theorem
\ref{thm:asym}'s main term; the window endpoints in Table \ref{tab:window} are the two roots of
the right-hand side, the lower at $t_{-}\approx(1-B/A)/(2\omega)<0.0017\,\omega$.

For the endpoint law, note that the two branches of $\nu$ cross at $L_x=8928.2$, and that for
$L\le L_x$ one has $\nu(L)=1/(cL)$ with $c=5.558691$. By the Laplace method, the exponent
$-L-t/(cL)$ is maximised at $L^{*}=\sqrt{t/c}$ with value $-2\sqrt{t/c}$; splitting off
$e^{-L/\sqrt t}$ to control the polynomial factor shows
$E(t)=\exp\bigl(-2\sqrt{t/c}\,(1+o(1))\bigr)$, valid while $L^{*}\le L_x$, i.e.\
$t\le cL_x^{2}\approx4.4\cdot10^{8}$ --- far beyond every entry of Table \ref{tab:window}.
Solving $(\tfrac12-\omega)t_{+}=2\sqrt{t_{+}/c}\,(1+o(1))$ gives
$t_{+}=(4/c+o(1))(\tfrac12-\omega)^{-2}$; numerically the law is accurate to $1\%$ already at
$\omega=0.499$ ($0.7196\cdot10^{6}$ predicted against $726\,493$ computed). Beyond the crossover
--- relevant only once $\tfrac12-\omega\lesssim2\cdot10^{-5}$ --- the Korobov--Vinogradov branch
takes over and the same computation gives $t_{+}\asymp(\tfrac12-\omega)^{-5/2+o(1)}$.
\end{proof}

\begin{remark}\label{rem:nonumerics}
Proposition \ref{prop:window} delimits what computation can say about the criterion. In the
variable $x=e^{t}$ of the underlying prime sums, the window reaches $x\approx10^{28}$ at
$\omega=0.1$, $x\approx10^{48}$ at $\omega=\tfrac14$, $x\approx10^{131}$ at $\omega=0.4$, and
$x\approx10^{3154}$ at $\omega=0.49$: any computation of $\Psi_\omega$ at feasible arguments takes
place deep inside a region where its positivity is already unconditionally known, and therefore
carries no information about $\Theta$ beyond the inputs it is built on. In particular, ``compute
$\Psi_\omega$ at large $t$ and watch for a dip'' is not a probe of $\Theta$. The two inputs act
separably: the verification height sets the \emph{level} of the obstruction bound, so the window
grows with $T_v$ like $t_{+}\sim\frac{2\log T_v}{1-2\omega}$ at moderate $\omega$, while the
zero-free region sets its \emph{decay}, which is invisible at moderate $\omega$ (a $1$--$7\%$
enlargement of $t_+$ for $\omega\le0.45$) and dominant at the PNT end, where it produces the
endpoint law $t_{+}=(4/c+o(1))(\tfrac12-\omega)^{-2}$ and reconnects the family continuously to
the unconditional case $\omega\ge\tfrac12$. This is the honest quantitative content available to
the criterion short of a genuine zero-free half-plane.
\end{remark}

% ===========================================================================
\section{The rate at which a violating zero is detected}\label{sec:sens}
% ===========================================================================

\begin{proof}[Proof of Proposition \ref{prop:sens}]
By \eqref{eq:gammarho} the quartet $\{\rho_0,\overline{\rho_0},1-\rho_0,1-\overline{\rho_0}\}$
corresponds to $\gamma\in\{\pm\tau_0\pm i\delta\}$. Fix such a $\gamma=a+i\delta$ with
$a=\pm\tau_0$. In \eqref{eq:Rclean}, the term $e^{-i\gamma t}=e^{-iat}e^{\delta t}$ dominates its
partner $e^{i\gamma t}=e^{iat}e^{-\delta t}$ as $t\to\infty$, so the contribution of this $\gamma$
to the oscillatory part is
\[
  -\frac{e^{(\delta-\omega)t}e^{-iat}}{2(\gamma-i\omega)^{2}}\bigl(1+O(e^{-2\delta t})\bigr),
  \qquad \gamma-i\omega=a+i(\delta-\omega).
\]
Summing over $a=\pm\tau_0$ and adding the contribution from the two $\gamma$ with
$\Ip\gamma=-\delta$ (for which the dominant term is $e^{i\gamma t}$, with denominator
$(\gamma+i\omega)^{2}$) yields
\[
  -\,e^{(\delta-\omega)t}\,\Rp\!\left[\frac{e^{-i\tau_0 t}}{\bigl(\tau_0+i(\delta-\omega)\bigr)^{2}}
   +\frac{e^{i\tau_0 t}}{\bigl(\tau_0-i(\delta-\omega)\bigr)^{2}}\right](1+o(1))
  =-\,\frac{2\,e^{(\delta-\omega)t}}{\tau_0^{2}+(\delta-\omega)^{2}}\cos(\tau_0t+\alpha)(1+o(1)),
\]
where $\alpha=2\arg\bigl(\tau_0+i(\delta-\omega)\bigr)$. The non-oscillatory part of the quartet's
contribution is $\omega A_0t+B_0=O(t)$ with $A_0,B_0$ the corresponding partial sums in
\eqref{eq:AB}. The $\limsup$ statement follows since $\cos(\tau_0t+\alpha)$ attains $\pm1$ on a
relatively dense set.
\end{proof}

\begin{remark}\label{rem:sensitivity}
Proposition \ref{prop:sens} should be contrasted with the Nyman--Beurling criterion, in which RH is
equivalent to $d_N\to0$ for an explicit sequence of Hilbert-space distances with
$\liminf d_N^{2}\log N\ge\sum_\rho m_\rho^{2}|\rho|^{-2}$ \cite{BDBLS,BurnolLower}. There, by an
identity of Balazard--Saias--Yor \cite{BSY}, the natural size of the violation signal of a zero
$\rho=\tfrac12+\delta+i\tau$ is $\log|\rho/(1-\rho)|\approx\delta/(\tfrac14+\tau^{2})$, to be seen
against a baseline of order $1/\log N$: detection requires $\log N\gtrsim\tau^{2}/\delta$, i.e.\
$N\ge\exp(c\,\tau^{2}/\delta)$. Here, in the same variable $t=\log x$, the same zero drives
$\Psi_\omega$ with an envelope growing like $e^{(\delta-\omega)t}$ and explicit amplitude:
combining Proposition \ref{prop:sens} with Corollary \ref{cor:margin}, the height at which a
violating quartet at fixed $\tau_0$ overtakes the margin is
$t\asymp(\delta-\omega)^{-1}\log\bigl((\delta-\omega)^{-1}\bigr)$. The criterion is exponentially
sensitive where Nyman--Beurling is nearly blind --- though Proposition \ref{prop:window} shows this
sensitivity, too, begins only beyond every feasible argument, since a violating zero must lie above
the verification height.
\end{remark}

% ===========================================================================
\section{Dirichlet $L$-functions, and the extended Selberg class}\label{sec:dirichlet}
% ===========================================================================

Everything so far used four inputs: the Hadamard expansion, the functional-equation antisymmetry
of $\xi'/\xi$ about $s=\tfrac12$, the strip bound $|\Ip\gamma|\le\tfrac12$, and the absence of
real zeros in $(\tfrac12,1]$. All four are available for real primitive Dirichlet characters, so
the theory transfers whole. What changes is instructive: the margin grows with the conductor, and
the absence-of-real-zeros input --- free for $\zeta$ --- becomes both necessary and visible.

Let $\chi$ be a real primitive character mod $q$, $a=(1-\chi(-1))/2$, and
\[
  \xi(s,\chi)=(q/\pi)^{(s+a)/2}\,\Gamma\!\left(\tfrac{s+a}2\right)L(s,\chi),
\]
entire of order one with $\xi(s,\chi)=\xi(1-s,\chi)$, since the root number of a real primitive
character equals $+1$ by Gauss's theorem (see \cite[\S9]{Davenport}). Because the coefficients are
real, zeros also pair under conjugation, so the zero set of $\Xi_\chi(z)=\xi(\tfrac12-iz,\chi)$
carries the full symmetry $\gamma\mapsto-\gamma$, $\gamma\mapsto\overline\gamma$ used throughout,
and $\Psi_{\omega,\chi}$ is defined by the series \eqref{eq:Psiomegaseries} over these $\gamma$.
We restrict to $q\le4\cdot10^{5}$, where the numerical verification of GRH by Platt
\cite{Platt2016} guarantees that $L(s,\chi)$ has no real zeros in $(0,1)$ and no non-real zeros of
height below $\max(10^{8}/q,200)$; this replaces the two numerical inputs used for $\zeta$.

\begin{theorem}\label{thm:chi}
Let $\chi$ be real primitive of modulus $q\le4\cdot10^{5}$ and $0<\omega<\tfrac12$, with sums over
the zeros $\gamma$ of $\Xi_\chi$.
\begin{enumerate}
\item[(i)] Unconditionally,
$A_\chi(\omega):=\sum_\gamma(\gamma^{2}+\omega^{2})^{-1}
 =\omega^{-1}(\xi'/\xi)(\tfrac12+\omega,\chi)$ and
$B_\chi(\omega):=\sum_\gamma(\gamma^{2}-\omega^{2})(\gamma^{2}+\omega^{2})^{-2}
 =(\xi'/\xi)'(\tfrac12+\omega,\chi)$; moreover $(\xi'/\xi)(\tfrac12,\chi)=0$, so the margin slope
vanishes at $\omega=0$ exactly as for $\zeta$.
\item[(ii)] If $\vartheta_\chi:=\sup_\rho\Rp(\rho)-\tfrac12\le\omega$ then
\[
  \Psi_{\omega,\chi}(t)=\lxi{\left(\tfrac12+\omega,\chi\right)}t
  +\lxi{'\left(\tfrac12+\omega,\chi\right)}
  +O\!\left(A^{*}_\chi(\omega)\,e^{-(\omega-\vartheta_\chi)t}\right),
\]
and $\Psi_{\omega,\chi}(t)>0$ for every
$t>(A^{*}_\chi(\omega)-B_\chi(\omega))/(\omega A_\chi(\omega))$, with
$A^{*}_\chi-B_\chi\le8\omega^{2}\sum_\rho(\Ip\rho)^{-4}$.
\item[(iii)] The analogue of Theorem \ref{thm:transition} holds verbatim.
\end{enumerate}
\end{theorem}

\begin{proof}
The proofs of Sections \ref{sec:setup}--\ref{sec:transition} apply with $\xi(\cdot)$ replaced by
$\xi(\cdot,\chi)$. Oddness of $(\xi'/\xi)(\cdot,\chi)$ about $\tfrac12$ follows from the even
functional equation, which gives $(\xi'/\xi)(\tfrac12,\chi)=0$ and removes the constant from the
Hadamard expansion of $\Xi_\chi'/\Xi_\chi$, exactly as in \cite[\S2.2]{Suzuki2023}. The strip
bound $|\Ip\gamma|\le\tfrac12$ holds as for $\zeta$. By \cite{Platt2016} there are no real zeros
in $(0,1)$ (so no purely imaginary $\gamma$, and no poles of $(\xi'/\xi)(\cdot,\chi)$ on
$(\tfrac12,1]$, which part (iii) needs), and every non-real zero off the critical line has height
above $200$, so $\Rp(\gamma^{2}+\omega^{2})>0$ for every $\gamma$. The pole terms $1/s+1/(s-1)$
are absent throughout, since $\xi(s,\chi)$ is entire.
\end{proof}

Table \ref{tab:conductor} records the margin across conductors: $B_\chi(0)=(\xi'/\xi)'(\tfrac12,\chi)$
--- the analogue of Remark \ref{rem:B0} --- and the slope at $\omega=\tfrac14$. The margin grows
with the conductor, because the lowest zero descends: \emph{within this family, $\zeta$ is the most
marginal member}.

\begin{table}[htbp]
\centering
\begin{tabular}{@{}lcc@{}}
\toprule
$L$-function & $B(0)$ & slope $(\xi'/\xi)(\tfrac34,\cdot)$ \\
\midrule
$\zeta$        & $0.046210$ & $0.011552$\\
$\chi$ mod $3$ & $0.113402$ & $0.028340$\\
$\chi$ mod $4$ & $0.156030$ & $0.038978$\\
$\chi$ mod $5$ & $0.156914$ & $0.039206$\\
\bottomrule
\end{tabular}
\vspace{2mm}
\caption{The margin across conductors.}
\label{tab:conductor}
\end{table}

The arithmetic side mirrors \eqref{eq:Psiomegaprimes}, with the pole block absent.

\begin{proposition}\label{prop:chiprime}
Let $\chi,q,a$ be as above and $\alpha:=(\tfrac12+\omega+a)/2$. For $t>0$,
\[
  \Psi_{\omega,\chi}(t)=\frac t2\left[\psi(\alpha)+\log\frac q\pi\right]
  +\frac14\left[\psi^{(1)}(\alpha)-e^{-2\alpha t}\,\Phi\!\left(e^{-2t},2,\alpha\right)\right]
  -\sum_{n\le e^{t}}\frac{\Lambda(n)\chi(n)}{n^{\frac12+\omega}}\,(t-\log n).
\]
\end{proposition}

\begin{proof}
By uniqueness of Fourier transforms it suffices to show that the right-hand side has Laplace
transform $-z^{-2}(\xi'/\xi)(\tfrac12+\omega-iz,\chi)$ for $\Ip z$ large, since the series
defining $\Psi_{\omega,\chi}$ has this transform by the argument of
\cite[\S\S2.2, 11]{Suzuki2023}, using $(\xi'/\xi)(\tfrac12,\chi)=0$. Write $s=\tfrac12+\omega-iz$;
for $\Rp s>1$,
\[
  \lxi{(s,\chi)}=\tfrac12\log\frac q\pi+\tfrac12\,\psi\!\left(\tfrac{s+a}2\right)
  -\sum_{n\ge2}\Lambda(n)\chi(n)\,n^{-s},
\]
with no pole terms. The transforms $\int_0^\infty te^{izt}dt=-z^{-2}$ and
$\int_0^\infty(t-\log n)\,n^{-\frac12-\omega}\mathbf 1_{[\log n,\infty)}(t)\,e^{izt}dt
 =-z^{-2}n^{-s}$ handle the first and third pieces. For the middle piece we claim, for $\alpha>0$
and $\Ip z>0$,
\begin{equation}\label{eq:gammalemma}
  \int_0^\infty\frac14\left[\psi^{(1)}(\alpha)-e^{-2\alpha t}\Phi(e^{-2t},2,\alpha)\right]
  e^{izt}\,dt
  =-\frac1{2z^{2}}\left[\psi\!\left(\alpha-\tfrac{iz}2\right)-\psi(\alpha)\right].
\end{equation}
Indeed $\psi^{(1)}(\alpha)-e^{-2\alpha t}\Phi(e^{-2t},2,\alpha)
 =\sum_{k\ge0}(k+\alpha)^{-2}\bigl(1-e^{-2(k+\alpha)t}\bigr)$, and per $k$, with
$c=2(k+\alpha)$, $\int_0^\infty(1-e^{-ct})e^{izt}dt=\tfrac iz-\tfrac1{c-iz}$; the identity
\[
  \frac1{4(k+\alpha)^{2}}\left[\frac iz-\frac1{2(k+\alpha)-iz}\right]
  =-\frac1{2z^{2}}\left[\frac1{k+\alpha}-\frac1{k+\alpha-iz/2}\right]
\]
is elementary --- both sides equal
$\tfrac{i}{4z}\bigl[(k+\alpha)(k+\alpha-\tfrac{iz}2)\bigr]^{-1}$ --- and summing over $k$ gives
\eqref{eq:gammalemma}. Since $(s+a)/2=\alpha-iz/2$, assembling the three pieces yields
$-z^{-2}(\xi'/\xi)(s,\chi)$ for $\Ip z>\tfrac12$, hence for $\Ip z>\tfrac12-\omega$ by analytic
continuation. Finally $\Psi_{\omega,\chi}(0)=0$, since
$\Phi(1,2,\alpha)=\zeta(2,\alpha)=\psi^{(1)}(\alpha)$.
\end{proof}

Table \ref{tab:chiprime} tests Proposition \ref{prop:chiprime} against the main term of Theorem
\ref{thm:chi}(ii), using prime powers $p^{m}\le e^{15}$ and no information about the zeros; the
difference sits inside the envelope $A_\chi(\omega)e^{-\omega t}$ in every row and decays at the
predicted rate. (The differences are roughly ten times the $\zeta$ case, as they must be:
$A_\chi$ is about three times larger and the lowest zeros are lower --- $6.02$ for $q=4$ against
$14.13$ --- so the leading damped oscillation is bigger.)

\begin{table}[htbp]
\centering
\begin{tabular}{@{}lllcccc@{}}
\toprule
$q$ & $\omega$ & $t$ & $\Psi_{\omega,\chi}(t)$ (primes) & main term & difference & $A_\chi(\omega)e^{-\omega t}$\\
\midrule
$3$ & $0.2$ & $6$  & $0.25378628$ & $0.24936876$ & $4.4\times10^{-3}$  & $3.4\times10^{-2}$\\
$3$ & $0.2$ & $15$ & $0.45339439$ & $0.45344292$ & $-4.9\times10^{-5}$ & $5.6\times10^{-3}$\\
$3$ & $0.3$ & $15$ & $0.62323052$ & $0.62324617$ & $-1.6\times10^{-5}$ & $1.3\times10^{-3}$\\
$3$ & $0.4$ & $15$ & $0.79281810$ & $0.79282263$ & $-4.5\times10^{-6}$ & $2.8\times10^{-4}$\\
$4$ & $0.2$ & $6$  & $0.35287485$ & $0.34295358$ & $9.9\times10^{-3}$  & $4.7\times10^{-2}$\\
$4$ & $0.2$ & $15$ & $0.62564343$ & $0.62367332$ & $2.0\times10^{-3}$  & $7.8\times10^{-3}$\\
$4$ & $0.3$ & $15$ & $0.85733527$ & $0.85691254$ & $4.2\times10^{-4}$  & $1.7\times10^{-3}$\\
$4$ & $0.4$ & $15$ & $1.08963532$ & $1.08954490$ & $9.0\times10^{-5}$  & $3.9\times10^{-4}$\\
$5$ & $0.2$ & $6$  & $0.36369455$ & $0.34497036$ & $1.9\times10^{-2}$  & $4.7\times10^{-2}$\\
$5$ & $0.2$ & $15$ & $0.62696177$ & $0.62731266$ & $-3.5\times10^{-4}$ & $7.8\times10^{-3}$\\
$5$ & $0.3$ & $15$ & $0.86199461$ & $0.86206119$ & $-6.7\times10^{-5}$ & $1.7\times10^{-3}$\\
$5$ & $0.4$ & $15$ & $1.09632744$ & $1.09633964$ & $-1.2\times10^{-5}$ & $3.9\times10^{-4}$\\
\bottomrule
\end{tabular}
\vspace{2mm}
\caption{Proposition \ref{prop:chiprime} tested against the arithmetic side. No zeros are used.}
\label{tab:chiprime}
\end{table}

The genuinely new phenomenon is what the criterion \emph{cannot} see. A real zero enters the
series with a locked phase, and pushes $\Psi_{\omega,\chi}$ up.

\begin{proposition}[Blindness to real zeros]\label{prop:siegel}
Let $\chi$ be a real primitive character, $0<\omega<\tfrac12$, and assume $\tfrac12+\omega$ is not
a zero of $L(\cdot,\chi)$. Write $\beta_0\in[\tfrac12,1)$ for the largest real zero of
$L(\cdot,\chi)$ in $[\tfrac12,1)$, with $\beta_0=\tfrac12$ if there is none, and
$\vartheta_{\mathrm{nr}}:=\sup\{\Rp\rho-\tfrac12:\rho\ \text{non-real}\}$.
\begin{enumerate}
\item[(i)] A real zero pair $\{\tfrac12+\delta,\tfrac12-\delta\}$ with $\delta>\omega$ contributes
\[
  \frac{e^{(\delta-\omega)t}}{(\delta-\omega)^{2}}
  +\frac{e^{-(\delta+\omega)t}}{(\delta+\omega)^{2}}\;>\;0
\]
to the oscillatory part \eqref{eq:Rclean} of $\Psi_{\omega,\chi}$.
\item[(ii)] If $\Psi_{\omega,\chi}(t)\ge0$ for all sufficiently large $t$, then
$\vartheta_{\mathrm{nr}}\le\max\bigl(\omega,\ \beta_0-\tfrac12\bigr)$.
\item[(iii)] Conversely, suppose $\vartheta_{\mathrm{nr}}\le\max(\omega,\beta_0-\tfrac12)$, that
every non-real zero satisfies $|\Ip\rho|\ge\tfrac12$, and, in the case $\beta_0>\tfrac12+\omega$,
that $2\sum_{\rho\ \mathrm{non\text{-}real}}(\Ip\rho)^{-2}<(\beta_0-\tfrac12-\omega)^{-2}$. Then
$\Psi_{\omega,\chi}(t)\ge0$ for all sufficiently large $t$.
\end{enumerate}
For $q\le4\cdot10^{5}$ the side conditions in (iii) hold by \cite{Platt2016} (indeed there are no
real zeros at all), so (ii) and (iii) are an equivalence. In particular the shifted criterion
cannot detect Siegel zeros, and the real-zero hypothesis in \cite[Thm.~1.1]{SuzukiSelberg} cannot
be dropped.
\end{proposition}

\begin{proof}
(i) For $\gamma=\pm i\delta$ one has $e^{\mp i\gamma t}=e^{\pm\delta t}$ and
$(\gamma\mp i\omega)^{2}=-(\delta\mp\omega)^{2}$, so each of the four terms that the pair
contributes to \eqref{eq:Rclean} is positive; summing them gives the display.

(ii) This is the Landau argument of Theorem \ref{thm:transition}(ii). Suppose
$\Psi_{\omega,\chi}\ge0$ eventually; its Laplace transform $F$ converges in a half-plane and
equals $-z^{-2}(\xi'/\xi)(\tfrac12+\omega-iz,\chi)$ plus an entire correction. By Landau's
theorem, $i\sigma_c$ is a singularity of $F$; the only candidates on the closed positive imaginary
axis are $z=0$ and points $i(\beta-\tfrac12-\omega)$ with $\beta$ a real zero, so
$\sigma_c\le\max(0,\beta_0-\tfrac12-\omega)$. Since $F$ is analytic above $\sigma_c$ and agrees
with the meromorphic expression, every zero satisfies
$\Rp\rho-\tfrac12-\omega\le\sigma_c$, which is the claim.

(iii) Write $\Psi_{\omega,\chi}=\omega A_\chi t+B_\chi+R$ as in \eqref{eq:Rclean}. If
$\beta_0\le\tfrac12+\omega$, all zeros satisfy $\Rp\rho\le\tfrac12+\omega$: real-zero pairs with
$\delta\le\omega$ contribute positively to $A_\chi$ (their terms are $(\omega^{2}-\delta^{2})^{-1}>0$),
non-real zeros contribute $\Rp[(\gamma^{2}+\omega^{2})^{-1}]>0$ by $|\Ip\rho|\ge\tfrac12\ge|b|$, so
$A_\chi(\omega)>0$, $|R|\le A^{*}_\chi<\infty$, and the margin wins for large $t$. If
$\beta_0=\tfrac12+\delta_r>\tfrac12+\omega$, then by (i) the largest real pair contributes at least
$e^{(\delta_r-\omega)t}(\delta_r-\omega)^{-2}(1+o(1))$, all other real pairs contribute positively,
and the non-real zeros contribute at most
$e^{(\delta_r-\omega)t}\cdot2\sum_{\mathrm{nr}}(\Ip\rho)^{-2}$ in absolute value, which the
amplitude condition makes strictly smaller; the linear terms are $O(t)$. Hence
$\Psi_{\omega,\chi}\to+\infty$.
\end{proof}

\begin{remark}[The extended Selberg class]\label{rem:selberg}
Suzuki has extended the unshifted theory to the extended Selberg class $S^{\sharp}$
\cite{SuzukiSelberg}: the screw function there acquires the extra terms
$-iB_Ft-\tfrac{m_0}2t^{2}$ (with $iB_F=(\xi_F'/\xi_F)(\tfrac12)$ and $m_0$ the order of the
central zero), the criterion concerns $\Rp(-g_F)$, is stated under the hypothesis that $F$ has no
real zeros off $s=\tfrac12$, and takes the ``eventually non-negative'' form even without a shift
--- which makes an explicit threshold in the style of Theorem \ref{thm:threshold} more, not less,
relevant there. Three remarks on how the present results sit in that generality. First, the
functional equation $\xi_F(s)=\omega_F\overline{\xi_F(1-\bar s)}$ forces
$\Rp(\xi_F'/\xi_F)(\tfrac12)=0$ for \emph{every} $F\in S^{\sharp}$, so the margin slope
$\Rp(\xi_F'/\xi_F)(\tfrac12+\omega)$ vanishes at $\omega=0$ always: the marginality mechanism of
Corollary \ref{cor:margin} is universal, the non-self-dual case being handled through the pair
$(F,\bar F)$, which restores the $\pm\gamma$ symmetry used in the proofs. Second, a central zero
of order $m_0$ contributes the non-negative term $m_0\,[\,t/\omega-(1-e^{-\omega t})/\omega^{2}]$
to the shifted function, so families with forced central zeros are non-marginal at every
$\omega>0$, with slope blowing up like $m_0/\omega$ as $\omega\to0^{+}$. Third, Proposition
\ref{prop:siegel} is the quantitative reason for the real-zero hypothesis in
\cite[Thm.~1.1]{SuzukiSelberg}: real zeros push the screw function up, at the fastest exponential
rate present, and no criterion of this family can see them.
\end{remark}

% ===========================================================================
\section{Numerical verification}\label{sec:numerics}
% ===========================================================================

All computations were done in \texttt{mpmath} at $20$--$30$ decimal digits, with $\Lambda(n)$
obtained by sieving. Scripts are available from the author. The Dirichlet computations ---
functional-equation and reality checks for $\xi(\cdot,\chi)$, zero counts against the
Riemann--von Mangoldt formula, the closed forms of Theorem \ref{thm:chi}(i) against direct zero
sums, and the arithmetic side --- are reported in Section \ref{sec:dirichlet}
(Tables \ref{tab:conductor} and \ref{tab:chiprime}).

\subsection{The arithmetic side}
As a check on the implementation of \eqref{eq:Psiomegaprimes} we reproduce a value printed by
Suzuki. In the proof of \cite[Thm.~4.1]{Suzuki2023} he records $\Psi(t_2)=0.0396618\ldots$ at
$t_2=0.464002\ldots$. Evaluating \eqref{eq:Psiomegaprimes} at $\omega=0$ gives
\[
  \Psi(0.464002)=0.03966175656\ldots,\qquad \Psi(0)=0 .
\]
Separately, the per-zero identity between the definition \eqref{eq:Psiomegadef} and the series
\eqref{eq:Psiomegaseries}, and the decomposition \eqref{eq:Rdecomp}, were verified symbolically and
to $30$ digits at complex $\gamma$.

\subsection{The closed forms of Proposition \ref{prop:AB}}
Table \ref{tab:AB} compares the closed forms against direct summation over the first $250$ zeros
together with a Riemann--von Mangoldt tail estimate $dN=(2\pi)^{-1}\log(\tau/2\pi)\,d\tau$. The
relative discrepancy is $1.13\times10^{-4}$ in every row, independent of $\omega$; this is the error
of the tail estimate, not of the identity.

\begin{table}[htbp]
\centering
\begin{tabular}{@{}lcccc@{}}
\toprule
$\omega$ & $A(\omega)$ (sum + tail) & $A(\omega)$ (closed form) & $B(\omega)$ (sum + tail) & $B(\omega)$ (closed form)\\
\midrule
$0.10$ & $0.04621447$ & $0.04620924$ & $0.04621299$ & $0.04620776$\\
$0.25$ & $0.04621057$ & $0.04620534$ & $0.04620128$ & $0.04619605$\\
$0.40$ & $0.04620333$ & $0.04619810$ & $0.04617957$ & $0.04617434$\\
\bottomrule
\end{tabular}
\vspace{2mm}
\caption{Proposition \ref{prop:AB} against direct summation over $250$ zeros plus tail.}
\label{tab:AB}
\end{table}

\subsection{The asymptotic, tested against the primes}
Table \ref{tab:asym} is the substantive check: the left column is computed from
\eqref{eq:Psiomegaprimes}, i.e.\ from prime powers $p^m\le e^{15}$ and elementary functions, using
no information about the zeros whatsoever; the right column is the main term of Theorem
\ref{thm:asym}. The difference lies inside the envelope $A(\omega)e^{-\omega t}$ of Theorem
\ref{thm:asym} in every row and decays at the predicted rate.

\begin{table}[htbp]
\centering
\begin{tabular}{@{}llcccc@{}}
\toprule
$\omega$ & $t$ & $\Psi_\omega(t)$ from \eqref{eq:Psiomegaprimes} & main term & difference & $A(\omega)e^{-\omega t}$\\
\midrule
$0.20$ & $6$  & $0.10222000$ & $0.10164948$ & $5.7\times10^{-4}$ & $1.4\times10^{-2}$\\
$0.20$ & $15$ & $0.18501313$ & $0.18482211$ & $1.9\times10^{-4}$ & $2.3\times10^{-3}$\\
$0.30$ & $6$  & $0.12966687$ & $0.12935586$ & $3.1\times10^{-4}$ & $7.6\times10^{-3}$\\
$0.30$ & $15$ & $0.25414919$ & $0.25410476$ & $4.4\times10^{-5}$ & $5.1\times10^{-4}$\\
$0.40$ & $6$  & $0.15721911$ & $0.15704977$ & $1.7\times10^{-4}$ & $4.2\times10^{-3}$\\
$0.40$ & $15$ & $0.32337324$ & $0.32336293$ & $1.0\times10^{-5}$ & $1.2\times10^{-4}$\\
\bottomrule
\end{tabular}
\vspace{2mm}
\caption{Theorem \ref{thm:asym} tested against the arithmetic side. No zeros are used.}
\label{tab:asym}
\end{table}

\subsection{The margin and the threshold}
Table \ref{tab:margin} records the slope, the intercept, the ratio $r:=B/A$, the threshold
$t_0=(A-B)/(\omega A)=(1-r)/\omega$ that Theorem \ref{thm:threshold} yields in the RH case
($A^{*}=A$), and the bound $8\omega^{2}\sum_\rho(\Ip\rho)^{-4}/(\xi'/\xi)(\tfrac12+\omega)$ valid
for every $\vartheta\le\omega$. The row $\omega=0$ returns $1.2\times10^{-30}$ for the slope, i.e.\
zero to working precision, confirming $(\xi'/\xi)(\tfrac12)=0$.

\begin{table}[htbp]
\centering
\begin{tabular}{@{}lccccc@{}}
\toprule
$\omega$ & $(\xi'/\xi)(\tfrac12+\omega)$ & $B(\omega)$ & $r(\omega)$ & $t_0$ (RH case) & $t_0$ bound (any $\vartheta\le\omega$)\\
\midrule
$0.49$  & $0.0226342$   & $0.0461565$ & $0.99922857$ & $0.00157$   & $0.00631$\\
$0.40$  & $0.0184792$   & $0.0461743$ & $0.99948567$ & $0.00129$   & $0.00515$\\
$0.30$  & $0.0138610$   & $0.0461899$ & $0.99971057$ & $0.00096$   & $0.00386$\\
$0.20$  & $0.0092414$   & $0.0462011$ & $0.99987132$ & $0.00064$   & $0.00257$\\
$0.10$  & $0.0046209$   & $0.0462078$ & $0.99996782$ & $0.00032$   & $0.00129$\\
$0.05$  & $0.0023105$   & $0.0462094$ & $0.99999196$ & $0.00016$   & $0.00064$\\
$0.01$  & $0.00046210$  & $0.0462100$ & $0.99999968$ & $3.2\times10^{-5}$ & $1.3\times10^{-4}$\\
$0.001$ & $0.000046210$ & $0.0462100$ & $0.9999999968$ & $3.2\times10^{-6}$ & $1.3\times10^{-5}$\\
$0$     & $1.2\times10^{-30}$ & $0.0462100$ & --- & --- & ---\\
\bottomrule
\end{tabular}
\vspace{2mm}
\caption{The margin of Theorem \ref{thm:asym} and the thresholds of Theorem \ref{thm:threshold}.
$1-r(\omega)\approx0.0032\,\omega^{2}$, reflecting
$A-B=2\omega^{2}\sum_\gamma(\gamma^{2}+\omega^{2})^{-2}$.}
\label{tab:margin}
\end{table}

\subsection{The unconditional window}
Table \ref{tab:window} records the window of Proposition \ref{prop:window}: the two roots of
$\omega A t+B-Ae^{-\omega t}-e^{(\frac12-\omega)t}E(t)$, with, for comparison, the endpoint that
the undamped bound $E\equiv S=2.96\cdot10^{-12}$ would give. The zero-free damping is a
$1$--$7\%$ effect for $\omega\le0.45$ and dominant near $\omega=\tfrac12$, where the endpoint law
$t_{+}\approx0.7196\,(\tfrac12-\omega)^{-2}$ takes over (predicted $7.20\cdot10^{5}$ against
computed $726\,493$ at $\omega=0.499$).

\begin{table}[htbp]
\centering
\begin{tabular}{@{}lcccc@{}}
\toprule
$\omega$ & $t_{-}(\omega)$ & $t_{+}$ ($E\equiv S$) & $t_{+}(\omega)$ & $x_{+}=e^{t_{+}}$\\
\midrule
$0.05$  & $8.0\times10^{-5}$ & $55.06$  & $55.84$   & $10^{24.3}$\\
$0.10$  & $1.6\times10^{-4}$ & $63.67$  & $64.68$   & $10^{28.1}$\\
$0.20$  & $3.2\times10^{-4}$ & $87.98$  & $89.86$   & $10^{39.0}$\\
$0.25$  & $4.0\times10^{-4}$ & $107.19$ & $109.94$  & $10^{47.8}$\\
$0.30$  & $4.8\times10^{-4}$ & $136.02$ & $140.42$  & $10^{61.0}$\\
$0.40$  & $6.4\times10^{-4}$ & $282.06$ & $300.89$  & $10^{130.7}$\\
$0.45$  & $7.2\times10^{-4}$ & $580.82$ & $663.59$  & $10^{288.2}$\\
$0.49$  & $7.9\times10^{-4}$ & $3079.1$ & $7262.3$  & $10^{3154}$\\
$0.494$ & $7.9\times10^{-4}$ & ---      & $20\,178$  & $10^{8763}$\\
$0.497$ & $8.0\times10^{-4}$ & ---      & $80\,694$  & $10^{35\,050}$\\
$0.499$ & $8.0\times10^{-4}$ & ---      & $726\,493$ & $10^{315\,500}$\\
\bottomrule
\end{tabular}
\vspace{2mm}
\caption{Proposition \ref{prop:window}: unconditional positivity of $\Psi_\omega$ on
$(t_{-},t_{+})$, from RH verified to $3\cdot10^{12}$ \cite{PlattTrudgian} and the zero-free
regions of \cite{MTY}.}
\label{tab:window}
\end{table}

\subsection{Sensitivity}
Table \ref{tab:sens} measures the growth rate of the envelope of a synthetic violating quartet at
ordinate $\tau_0=50$, taking the maximum of $|\Psi_\omega^{(\rho_0)}|$ over one period of
$\cos(\tau_0t)$. Where the exponential has fully separated from the linear drift the measured rate
agrees with $\delta-\omega$ to five digits; where it has not, the drift contributes and the measured
rate is smaller. Rows with $\delta<\omega$ grow at the rate of the linear drift
$\omega A_0t+B_0$ alone: the measured rates $0.01833$ and $0.01638$ match
$\tfrac1{20}\log\bigl((60\,\omega A_0+B_0)/(40\,\omega A_0+B_0)\bigr)=0.01838$ and $0.01682$ to
within the residual damped oscillation, i.e.\ no exponential component survives, as predicted. Comparison against the amplitude
$2(\tau_0^{2}+(\delta-\omega)^{2})^{-1}e^{(\delta-\omega)t}$ of Proposition \ref{prop:sens} gives
agreement to $1$--$6\%$ at the heights tested, the discrepancy being the $O(t)$ term.

\begin{table}[htbp]
\centering
\begin{tabular}{@{}llccrr@{}}
\toprule
$\delta$ & $\omega$ & envelope at $t=40$ & envelope at $t=60$ & measured rate & $\delta-\omega$\\
\midrule
$0.30$ & $0.10$ & $2.417$               & $1.3195\times10^{2}$ & $+0.20000$ & $+0.2000$\\
$0.20$ & $0.10$ & $5.175\times10^{-2}$  & $3.341\times10^{-1}$ & $+0.09325$ & $+0.1000$\\
$0.15$ & $0.10$ & $1.392\times10^{-2}$  & $2.727\times10^{-2}$ & $+0.03364$ & $+0.0500$\\
$0.10$ & $0.20$ & $1.445\times10^{-2}$  & $2.084\times10^{-2}$ & $+0.01833$ & $-0.1000$\\
$0.05$ & $0.10$ & $8.114\times10^{-3}$  & $1.126\times10^{-2}$ & $+0.01638$ & $-0.0500$\\
\bottomrule
\end{tabular}
\vspace{2mm}
\caption{Proposition \ref{prop:sens}: envelope growth for a synthetic quartet at $\tau_0=50$.}
\label{tab:sens}
\end{table}

% ===========================================================================
\section{What this supplies, and what it does not}\label{sec:discussion}
% ===========================================================================

\subsection{The shape of the problem}
Theorem \ref{thm:asym} exhibits the criterion as a competition:
\begin{equation}\label{eq:competition}
  \Psi_\omega(t)=\underbrace{\lxi{\left(\tfrac12+\omega\right)}t+\lxi{'\left(\tfrac12+\omega\right)}}
                 _{\text{margin, linear in }t,\ \text{unconditional}}
  \;-\;\underbrace{e^{-\omega t}\times(\text{oscillation over zeros})}_{\text{obstruction}} .
\end{equation}
Under Hypothesis \ref{hyp} with $\vartheta\le\omega$ the obstruction is bounded by
$A^{*}(\omega)e^{-(\omega-\vartheta)t}$ and the margin wins beyond an explicit
$t_0(\omega)\le0.0129\,\omega$. If instead there is a zero at $\Rp\rho=\tfrac12+\delta$ with
$\delta>\omega$, the obstruction contains a term of size $e^{(\delta-\omega)t}$ and eventually
swamps the margin. Theorem \ref{thm:suzuki} is exactly this dichotomy; Theorems
\ref{thm:asym}--\ref{thm:threshold} and Proposition \ref{prop:sens} make both sides of it
quantitative.

\subsection{No free lunch}
To prove $\Psi_\omega\ge0$ for a fixed $\omega<\tfrac12$ from the arithmetic side one must show
that the fluctuation of
\[
  \sum_{n\le e^{t}}\frac{\Lambda(n)}{n^{\frac12+\omega}}(t-\log n)
\]
about its smooth part is bounded by the linear margin $(\xi'/\xi)(\tfrac12+\omega)\,t
\approx0.0462\,\omega\,t$. Without zero-free information that fluctuation may a priori be as large
as $e^{(\frac12-\omega)t}$. Converting an exponential a priori bound into a linear one is precisely
the content of a zero-free half-plane. The reformulation transports the difficulty faithfully; it
does not reduce it, and we make no claim that it does.

\subsection{Three respects in which the shifted case is nevertheless better behaved}

\emph{(i) The criterion is not saturated.} For $\omega=0$ the inequality $\Psi\ge0$ has
$\liminf\Psi=0$ conditionally \eqref{eq:liminf}, so no argument that establishes the inequality up to
a positive error can succeed. Corollary \ref{cor:margin} shows that this obstruction is a property
of the point $\omega=0$ alone: it is the vanishing of $\xi'/\xi$ at the centre of the critical
strip. For every $\omega>0$ there is a margin, it is explicit, and it grows without bound.

\emph{(ii) The required cancellation is linear, not square-root.} Most positivity criteria for RH
demand square-root cancellation in an arithmetic sum. Here one needs only $o(t)$ control of a
twice-integrated weighted Chebyshev error --- a much weaker requirement, bought by asking for a
weaker conclusion.

\emph{(iii) Failure is exponentially visible.} Proposition \ref{prop:sens} and Remark
\ref{rem:sensitivity}.

\subsection{Zero-density estimates: what they can and cannot do here}\label{sec:zerodensity}
The obstruction in \eqref{eq:competition} is, by \eqref{eq:Rclean} and \eqref{eq:gammarho}, a sum
over zeros weighted by $e^{(\Rp\rho-\frac12-\omega)t}$ and $|\rho|^{-2}$. Sums of this shape are
what zero-density theorems --- from Ingham to the recent improvement of Guth and Maynard
\cite{GuthMaynard} --- are built to control, so it is natural to ask whether, for some
$\omega<\tfrac12$, known density estimates inserted into \eqref{eq:Rclean} bound the obstruction by
the margin. The answer is no, for every $\omega$, and for a structural reason worth recording:
density estimates bound \emph{counts} $N(\sigma,T)$, and every known bound is consistent with the
existence of a single zero at any fixed $\beta<1$; a single permitted zero at
$\beta>\tfrac12+\omega$ contributes $e^{(\beta-\frac12-\omega)t}$ and defeats any linear margin as
$t\to\infty$. No argument using count upper bounds alone can close the gap --- to bound the
obstruction by the margin for all $t$ \emph{is} to assert the zero-free half-plane.

What density and zero-free-region technology does control is the \emph{onset}: it constrains where
in $t$ a violation could first appear. The Korobov--Vinogradov region forces a zero at
$\beta$ close to $1$ to have height $\tau\ge\exp\bigl(c(1-\beta)^{-3/2}\bigr)$, and a zero of
height $\tau$ enters \eqref{eq:Rclean} with weight $\tau^{-2}$; both effects postpone the possible
failure of $\Psi_\omega\ge0$ to explicit ranges of $t$, exactly as verification to height $T_v$
does in Proposition \ref{prop:window}. In this precise sense the criterion's open content lives
entirely in the tail $t\to\infty$, where only genuine zero-free-half-plane information can reach.
The productive question is therefore not whether density estimates prove the criterion --- they
cannot --- but how the two available inputs shape the unconditional window of Proposition
\ref{prop:window}, and there the answer is now quantitative. The verification height sets the
\emph{level} of the obstruction bound: $t_{+}\approx(\tfrac12-\omega)^{-1}\log(1/S)$ at moderate
$\omega$, so raising $T_v$ moves the window logarithmically. The zero-free region sets its
\emph{decay}: negligible at moderate $\omega$, it takes over at the PNT end and produces the
endpoint law $t_{+}=(4/c+o(1))(\tfrac12-\omega)^{-2}$, so the constant $c$ of \cite{MTY} enters
the criterion quadratically. All the action of zero-free technology on this family is
concentrated near $\omega=\tfrac12$.

\subsection{Status of the family}
We note finally that $\Psi_\omega$ appears to be unworked. As of August 2026, Suzuki's paper has
twelve recorded citations (Semantic Scholar), of which eight are his own; none of them takes up
$\Psi_\omega$ or Theorem \ref{thm:suzuki}. In particular the later work extending screw functions
to the extended Selberg class \cite{SuzukiSelberg} contains no shifted family --- the symbol
$\omega$ there is the root number --- and its ``eventually non-negative'' criterion, together
with its real-zero hypothesis, is exactly the point at which the shifted, quantitative theory of
the present paper attaches (Section \ref{sec:dirichlet}).

% ===========================================================================
\begin{thebibliography}{99}

\bibitem{BDBLS} L.~Báez-Duarte, M.~Balazard, B.~Landreau and E.~Saias,
\emph{Notes sur la fonction $\zeta$ de Riemann, 3}, Adv. Math. \textbf{149} (2000), 130--144.

\bibitem{BSY} M.~Balazard, E.~Saias and M.~Yor,
\emph{Notes sur la fonction $\zeta$ de Riemann, 2}, Adv. Math. \textbf{143} (1999), 284--287.

\bibitem{BombieriI} E.~Bombieri,
\emph{Remarks on Weil's quadratic functional in the theory of prime numbers, I},
Atti Accad. Naz. Lincei Cl. Sci. Fis. Mat. Natur. Rend. Lincei (9) Mat. Appl. \textbf{11} (2000),
183--233.

\bibitem{BombieriII} E.~Bombieri,
\emph{A variational approach to the explicit formula},
Comm. Pure Appl. Math. \textbf{56} (2003), 1151--1164.

\bibitem{BombieriLagarias} E.~Bombieri and J.~C.~Lagarias,
\emph{Complements to Li's criterion for the Riemann hypothesis},
J. Number Theory \textbf{77} (1999), 274--287.

\bibitem{BurnolLower} J.-F.~Burnol,
\emph{A lower bound in an approximation problem involving the zeros of the Riemann zeta function},
Adv. Math. \textbf{170} (2002), 56--70.

\bibitem{Davenport} H.~Davenport,
\emph{Multiplicative Number Theory}, 3rd ed., Graduate Texts in Mathematics \textbf{74},
Springer, New York, 2000.

\bibitem{Ford} K.~Ford,
\emph{Zero-free regions for the Riemann zeta function},
in: Number Theory for the Millennium II, A K Peters, 2002, 25--56.

\bibitem{Guinand} A.~P.~Guinand,
\emph{A summation formula in the theory of prime numbers},
Proc. London Math. Soc. (2) \textbf{50} (1949), 107--119.

\bibitem{GuthMaynard} L.~Guth and J.~Maynard,
\emph{New large value estimates for Dirichlet polynomials},
Ann. of Math. (2) \textbf{203} (2026), no.~2; arXiv:2405.20552.

\bibitem{KreinLanger} M.~G.~Kre\u\i n and H.~Langer,
\emph{Continuation of hermitian positive definite functions and related questions},
Integral Equations Operator Theory \textbf{78} (2014), 1--69.

\bibitem{Lagarias1999} J.~C.~Lagarias,
\emph{On a positivity property of the Riemann $\xi$-function},
Acta Arith. \textbf{89} (1999), 217--234.

\bibitem{Li} X.-J.~Li,
\emph{The positivity of a sequence of numbers and the Riemann hypothesis},
J. Number Theory \textbf{65} (1997), 325--333.

\bibitem{MTY} M.~J.~Mossinghoff, T.~S.~Trudgian and A.~Yang,
\emph{Explicit zero-free regions for the Riemann zeta-function},
Res. Number Theory \textbf{10} (2024), Paper No. 11; arXiv:2212.06867.

\bibitem{Platt2016} D.~J.~Platt,
\emph{Numerical computations concerning the GRH},
Math. Comp. \textbf{85} (2016), 3009--3027; arXiv:1305.3087.

\bibitem{PlattTrudgian} D.~Platt and T.~Trudgian,
\emph{The Riemann hypothesis is true up to $3\cdot10^{12}$},
Bull. Lond. Math. Soc. \textbf{53} (2021), 792--797.

\bibitem{Suzuki2023} M.~Suzuki,
\emph{Aspects of the screw function corresponding to the Riemann zeta-function},
J. London Math. Soc. (2) \textbf{108} (2023), 1448--1487; arXiv:2206.03682.

\bibitem{SuzukiSelberg} M.~Suzuki,
\emph{Screw functions of Dirichlet series in the extended Selberg class}, arXiv:2209.12832.

\bibitem{Titchmarsh} E.~C.~Titchmarsh,
\emph{The Theory of the Riemann Zeta-Function}, 2nd ed., revised by D.~R.~Heath-Brown,
Oxford Univ. Press, 1986.

\bibitem{TrudgianArg} T.~S.~Trudgian,
\emph{An improved upper bound for the argument of the Riemann zeta-function on the critical line
II}, J. Number Theory \textbf{134} (2014), 280--292.

\bibitem{Weil} A.~Weil,
\emph{Sur les ``formules explicites'' de la th\'eorie des nombres premiers},
Comm. S\'em. Math. Univ. Lund (1952), 252--265.

\bibitem{Widder} D.~V.~Widder,
\emph{The Laplace Transform}, Princeton Math. Series \textbf{6}, Princeton Univ. Press, 1941.

\bibitem{Yoshida} H.~Yoshida,
\emph{On Hermitian forms attached to zeta functions},
in: Zeta Functions in Geometry, Adv. Stud. Pure Math. \textbf{21}, Kinokuniya, 1992, 281--325.

\end{thebibliography}

\end{document}
